
\documentclass[11pt]{extarticle}

\usepackage{import}
\usepackage[english]{babel}
\usepackage{graphicx}
\usepackage{framed}
\usepackage[normalem]{ulem}
\usepackage{indentfirst}
\usepackage[T1]{fontenc}
%\usepackage{pifont} %For unusual symbols
%\usepackage{mathdots} %For unusual combinations of dots
\usepackage{wrapfig}
\usepackage[inline,shortlabels]{enumitem}
\setlist{topsep=2pt,itemsep=2pt,parsep=0pt,partopsep=0pt}
\usepackage[dvipsnames]{xcolor}
\usepackage[utf8]{inputenc}
\usepackage[a4paper, top=0.5in,bottom=0.2in, left=0.5in, right=0.5in, footskip=0.3in, includefoot]{geometry}
\usepackage[most]{tcolorbox}
\usepackage{tikz,tikz-3dplot,tikz-cd,tkz-tab,tkz-euclide,pgf,pgfplots}
\pgfplotsset{compat=newest}
\usepackage{multicol}
\usepackage[bottom,multiple]{footmisc} %ensures footnotes are at the bottom of the page, and separates footnotes by a comma if they are adjacent
\usepackage[backend=bibtex,style=alphabetic]{biblatex}
\renewcommand*{\finalnamedelim}{\addcomma\addspace} %forces authors' names to be separated by comma, instead of "and"
\addbibresource{bibliography}
\usepackage{hyperref}
\usepackage[nameinlink]{cleveref} %nameinlink ensures that the entire element is clickable in the pdf, not just the number


\usepackage{amsmath,amsthm,amssymb,amsfonts}
\usepackage[italicdiff]{physics}
\usepackage{lmodern,mathrsfs}
\usepackage{bbm}
\usepackage{bm}






% Math Special Letters
% \Letter means \matnbb{letter}
\newcommand{\1}{\mathbbm{1}}
\newcommand{\A}{\mathbb{A}}
\newcommand{\B}{\mathbb{B}}
\newcommand{\C}{\mathbb{C}}
\newcommand{\D}{\mathbb{D}}
\newcommand{\E}{\mathbb{E}}
\newcommand{\F}{\mathbb{F}}
\newcommand{\I}{\mathbb{I}}
\newcommand{\K}{\mathbb{K}}
\newcommand{\N}{\mathbb{N}}
\newcommand{\Q}{\mathbb{Q}}
\newcommand{\R}{\mathbb{R}}
\newcommand{\T}{\mathbb{T}}
\newcommand{\V}{\mathbb{V}}
\newcommand{\X}{\mathbb{X}}
\newcommand{\Y}{\mathbb{Y}}
\newcommand{\Z}{\mathbb{Z}}
\newcommand{\M}{\mathbb{M}}
\renewcommand{\P}{\mathbb{P}}

% \Letters means \mathcal{letter}
\newcommand{\As}{\mathcal{A}}
\newcommand{\Bs}{\mathcal{B}}
\newcommand{\Cs}{\mathcal{C}}
\newcommand{\Ds}{\mathcal{D}}
\newcommand{\Es}{\mathcal{E}}
\newcommand{\Fs}{\mathcal{F}}
\newcommand{\Gs}{\mathcal{G}}
\newcommand{\Hs}{\mathcal{H}}
\newcommand{\Is}{\mathcal{I}}
\newcommand{\Js}{\mathcal{J}}
\newcommand{\Ks}{\mathcal{K}}
\newcommand{\Ls}{\mathcal{L}}
\newcommand{\Ms}{\mathcal{M}}
\newcommand{\Ns}{\mathcal{N}}
\newcommand{\Os}{\mathcal{O}}
\newcommand{\Ps}{\mathcal{P}}
\newcommand{\Qs}{\mathcal{Q}}
\newcommand{\Rs}{\mathcal{R}}
\newcommand{\Ss}{\mathcal{S}}
\newcommand{\Ts}{\mathcal{T}}
\newcommand{\Us}{\mathcal{U}}
\newcommand{\Vs}{\mathcal{V}}
\newcommand{\Ws}{\mathcal{W}}
\newcommand{\Xs}{\mathcal{X}}
\newcommand{\Ys}{\mathcal{Y}}
\newcommand{\Zs}{\mathcal{Z}}

% \letterb means \textbf{letter}
\newcommand{\ab}{\textbf{a}}
\newcommand{\bb}{\textbf{b}}
\newcommand{\cb}{\textbf{c}}
\newcommand{\db}{\textbf{d}}
\newcommand{\ub}{\textbf{u}}
\newcommand{\wb}{\textbf{w}}
\newcommand{\xb}{\textbf{x}}
\newcommand{\yb}{\textbf{y}}
\newcommand{\zb}{\textbf{z}}
\newcommand{\Ab}{\textbf{A}}
\newcommand{\Bb}{\textbf{B}}
\newcommand{\Cb}{\textbf{C}}
\newcommand{\Db}{\textbf{D}}
\newcommand{\eb}{\textbf{e}}
\newcommand{\ex}{\textbf{e}_x}
\newcommand{\ey}{\textbf{e}_y}
\newcommand{\ez}{\textbf{e}_z}

% bar and hat
\renewcommand{\hat}[1]{\widehat{#1}}
\renewcommand{\bar}[1]{\overline{#1}}

% Some basic math operations
\newcommand{\inprod}[2]{\left\langle #1, #2 \right\rangle} % inner product
\renewcommand{\tr}[1]{\operatorname{tr}\left(#1\right)}	%trace
\newcommand{\floor}[1]{\ensuremath{\left\lfloor#1\right\rfloor}}  
\newcommand{\ceil}[1]{\ensuremath{\left\lceil#1\right\rceil}}

\newcommand{\limm}[2]{\displaystyle\lim_{#1 \to #2}}
\newcommand{\lsup}[2]{\displaystyle\limsup_{#1 \to #2}}
\newcommand{\linf}[2]{\displaystyle\liminf_{#1 \to #2}}
\newcommand{\summ}[2]{\sum\limits_{#1}^{#2}}


\renewcommand{\cap}{\bigcap}
\renewcommand{\cup}{\bigcup}
\newcommand{\kup}{\bigsqcup}

\newcommand{\capp}[2]{\bigcap_{#1}^{#2}}
\newcommand{\cupp}[2]{\bigcup_{#1}^{#2}}
\newcommand{\kupp}[2]{\bigsqcup_{#1}^{#2}}

\newcommand{\p}[3][]{\frac{\pd^{#1}#2}{\pd{#3}^{#1}}}	% partial derivative
\renewcommand{\d}{\mathrm{d}}                % differential
\newcommand{\proj}[2]{\operatorname{proj}_{#2}#1}			% projection
\newcommand{\vspan}[1]{\operatorname{span}\left(#1\right)}	% vector span
\newcommand{\diag}[1]{\operatorname{diag}\!\left(#1\right)}			% diagonal

\newcommand{\argmin}{\operatornamewithlimits{arg\,min}}
\newcommand{\argmax}{\operatornamewithlimits{arg\,max}}



% Math symbols
\newcommand{\vep}{\varepsilon}
\newcommand{\ep}{\epsilon}
\newcommand{\vphi}{\varphi}
\newcommand{\emp}{\emptyset}
\newcommand{\trs}{\intercal}

\newcommand{\dx}{\,dx}
\newcommand{\dy}{\,dy}
\newcommand{\dt}{\,dt}
\newcommand{\dr}[1]{\,d#1} %derivative

\newcommand{\<}{\langle}
\renewcommand{\>}{\rangle}

% Other easy controls
\newcommand{\todo}[1]{\textcolor{red}{\textbf{#1}}} %To remind me of unfinished work to fix later
\newcommand{\hide}[1]{} %To hide large blocks of code without using % symbols
\renewcommand{\iff}{\Leftrightarrow}
\newcommand{\np}{\vspace{5pt} \noindent}
\newcommand{\imply}{$\Rightarrow\ $}
\newcommand{\st}{\mathrm{s.t.}\ }


%Math Operators
\DeclareMathOperator{\range}{\text{range}}
\DeclareMathOperator{\domain}{\text{domain}}
\DeclareMathOperator{\im}{\text{im}}

\numberwithin{equation}{section}


\newtheoremstyle{mystyle}{}{}{}{}{\sffamily\bfseries}{.}{ }{}
\newtheoremstyle{cstyle}{}{}{}{}{\sffamily\bfseries}{.}{ }{\thmnote{#3}}
\makeatletter
\renewenvironment{proof}[1][\proofname] {\par\pushQED{\qed}{\normalfont\sffamily\bfseries\topsep6\p@\@plus6\p@\relax #1\@addpunct{.} }}{\popQED\endtrivlist\@endpefalse}
\makeatother
\newcommand{\coolqed}[1]{\includegraphics[width=#1cm]{sunglasses_emoji.png}} %Defines the new QED symbol
\renewcommand{\qedsymbol}{\coolqed{0.32}} %Implements the new QED symbol
\theoremstyle{mystyle}{\newtheorem{definition}{Definition}[section]}
\theoremstyle{mystyle}{\newtheorem{proposition}[definition]{Proposition}}
\theoremstyle{mystyle}{\newtheorem{theorem}[definition]{Theorem}}
\theoremstyle{mystyle}{\newtheorem{lemma}[definition]{Lemma}}
\theoremstyle{mystyle}{\newtheorem{corollary}[definition]{Corollary}}
\theoremstyle{mystyle}{\newtheorem*{claim}{Claim}}
\theoremstyle{mystyle}{\newtheorem*{remark}{Remark}}
\theoremstyle{mystyle}{\newtheorem*{remarks}{Remarks}}
\theoremstyle{mystyle}{\newtheorem*{example}{Example}}
\theoremstyle{mystyle}{\newtheorem*{examples}{Examples}}
\theoremstyle{definition}{\newtheorem*{exercise}{Exercise}}
\theoremstyle{cstyle}{\newtheorem*{cthm}{}}

%Warning environment
\newtheoremstyle{warn}{}{}{}{}{\normalfont}{}{ }{}
\theoremstyle{warn}
\newtheorem*{warning}{\warningsign{0.2}\relax}

%Symbol for the warning environment, designed to be easily scalable
\newcommand{\warningsign}[1]{\tikz[scale=#1,every node/.style={transform shape}]{\draw[-,line width={#1*0.8mm},red,fill=yellow,rounded corners={#1*2.5mm}] (0,0)--(1,{-sqrt(3)})--(-1,{-sqrt(3)})--cycle;
\node at (0,-1) {\fontsize{48}{60}\selectfont\bfseries!};}}

\tcolorboxenvironment{definition}{boxrule=0pt,boxsep=0pt,colback={red!10},left=8pt,right=8pt,enhanced jigsaw, borderline west={2pt}{0pt}{red},sharp corners,before skip=10pt,after skip=10pt,breakable}
\tcolorboxenvironment{proposition}{boxrule=0pt,boxsep=0pt,colback={Orange!10},left=8pt,right=8pt,enhanced jigsaw, borderline west={2pt}{0pt}{Orange},sharp corners,before skip=10pt,after skip=10pt,breakable}
\tcolorboxenvironment{theorem}{boxrule=0pt,boxsep=0pt,colback={blue!10},left=8pt,right=8pt,enhanced jigsaw, borderline west={2pt}{0pt}{blue},sharp corners,before skip=10pt,after skip=10pt,breakable}
\tcolorboxenvironment{lemma}{boxrule=0pt,boxsep=0pt,colback={Cyan!10},left=8pt,right=8pt,enhanced jigsaw, borderline west={2pt}{0pt}{Cyan},sharp corners,before skip=10pt,after skip=10pt,breakable}
\tcolorboxenvironment{corollary}{boxrule=0pt,boxsep=0pt,colback={violet!10},left=8pt,right=8pt,enhanced jigsaw, borderline west={2pt}{0pt}{violet},sharp corners,before skip=10pt,after skip=10pt,breakable}
\tcolorboxenvironment{proof}{boxrule=0pt,boxsep=0pt,blanker,borderline west={2pt}{0pt}{CadetBlue!80!white},left=8pt,right=8pt,sharp corners,before skip=10pt,after skip=10pt,breakable}
\tcolorboxenvironment{remark}{boxrule=0pt,boxsep=0pt,blanker,borderline west={2pt}{0pt}{Green},left=8pt,right=8pt,before skip=10pt,after skip=10pt,breakable}
\tcolorboxenvironment{remarks}{boxrule=0pt,boxsep=0pt,blanker,borderline west={2pt}{0pt}{Green},left=8pt,right=8pt,before skip=10pt,after skip=10pt,breakable}
\tcolorboxenvironment{example}{boxrule=0pt,boxsep=0pt,blanker,borderline west={2pt}{0pt}{Black},left=8pt,right=8pt,sharp corners,before skip=10pt,after skip=10pt,breakable}
\tcolorboxenvironment{examples}{boxrule=0pt,boxsep=0pt,blanker,borderline west={2pt}{0pt}{Black},left=8pt,right=8pt,sharp corners,before skip=10pt,after skip=10pt,breakable}
\tcolorboxenvironment{cthm}{boxrule=0pt,boxsep=0pt,colback={gray!10},left=8pt,right=8pt,enhanced jigsaw, borderline west={2pt}{0pt}{gray},sharp corners,before skip=10pt,after skip=10pt,breakable}
\tcolorboxenvironment{claim}{boxrule=0pt,boxsep=0pt,blanker,borderline west={2pt}{0pt}{Green},left=8pt,right=8pt,before skip=10pt,after skip=10pt,breakable}

%align and align* environments with inline size
\newenvironment{talign}{\let\displaystyle\textstyle\align}{\endalign}
\newenvironment{talign*}{\let\displaystyle\textstyle\csname align*\endcsname}{\endalign}

\usepackage[explicit]{titlesec}
\titleformat{\section}{\fontsize{16}{18}\sffamily\bfseries}{\thesection}{16pt}{#1}
\titleformat{\subsection}{\fontsize{12}{14}\sffamily\bfseries}{\thesubsection}{12pt}{#1}
\titleformat{\subsubsection}{\fontsize{10}{12}\sffamily\bfseries}{\thesubsubsection}{8pt}{#1}

\titlespacing*{\section}{0pt}{5pt}{5pt}
\titlespacing*{\subsection}{0pt}{5pt}{5pt}
\titlespacing*{\subsubsection}{0pt}{5pt}{5pt}

%\newcommand{\sectionbreak}{\clearpage} %Start every section on a new page

\newcommand{\Disp}{\displaystyle}
\newcommand{\qe}{\hfill\(\bigtriangledown\)}
\DeclareMathAlphabet\mathbfcal{OMS}{cmsy}{b}{n}
\setlength{\parindent}{0.2in}
\setlength{\parskip}{0pt}
\setlength{\columnseprule}{0pt}

\title{\huge\sffamily\bfseries MATH 602 Functional Analysis}
\author{\Large\sffamily Xinchang Wang}
\date{\sffamily FALL 2022 University of Michigan}




\begin{document}

\setlength{\abovedisplayskip}{3pt}
\setlength{\belowdisplayskip}{3pt}
\setlength{\abovedisplayshortskip}{0pt}
\setlength{\belowdisplayshortskip}{0pt}
\maketitle

%Custom colors for different environments
\definecolor{contcol1}{HTML}{72E094}
\definecolor{contcol2}{HTML}{24E2D6}
\definecolor{convcol1}{HTML}{C0392B}
\definecolor{convcol2}{HTML}{8E44AD}

\begin{tcolorbox}[title=Contents, fonttitle=\huge\sffamily\bfseries\selectfont,interior style={left color=contcol1!40!white,right color=contcol2!40!white},frame style={left color=contcol1!80!white,right color=contcol2!80!white},coltitle=black,top=2mm,bottom=2mm,left=2mm,right=2mm,drop fuzzy shadow,enhanced,breakable]
\makeatletter
\@starttoc{toc}
\makeatother
\end{tcolorbox}

\clearpage
\section{Introduction}
\subsection{Vector Space}
\begin{definition}[Linear Vector Space]\ \\
A vector space $V$ over a field $\F$ is a set closed under operations of addition and scalar multiplication satisfying the following rules ($u,v,w \in V,\ a,b \in \F$):
\begin{itemize}
    \item $u + (v+w) = (u+v)+w$.
    \item $u + v = v + u$.
    \item $\exists\ 0 \in V$ such that $ 0 + v = v$, $\forall\ v \in V$.
    \item $\forall\ v \in V$, $\exists\ -v \in V$ s.t. $v + (-v) = 0$.
    \item $a(bv) = (ab)v$.
    \item $1 v = v$, where $1$ denotes the \textit{multiplicative identity} in $F$.
    \item $a(u+v) = au + av$.
    \item $(a+b)v = av + bv$.
\end{itemize}
In this context, the elements of $V$ are commonly called vectors, and the elements of $\F$ are called scalars. If $\F = \R$, we call $V$ a \textit{Real Vector Space}, and if $\F = \C$, we call $V$ a \textit{Complex Vector Space}.
\end{definition}
\begin{example}
$\R^n (\C^n)$ is an $n$-dimensional real (complex) vector space.
\end{example}

\np In this course we mainly concentrate on $\infty$-dimensional vector space.

\begin{example}\ \\
$K$ is a compact Hausdorff space, $E = \{f: K \to \R: f(\cdot) \text{ is continuous }\}$ is an $\infty$ dim real vector space.
\end{example}


\vspace{6pt}
\subsection{Quotient Space}
\noindent Vector space can have many subspaces. Let $E = $ Vector Space, $E_1 \subset E$ a \underline{proper subspace} (i.e. $E_1 \neq E$).

\begin{definition}[Quotient Space]\ \\
The \textit{Quotient Space} $E / E_1$ is the set of \underline{equivalent classes} of vectors in $E$, where equivalence relation is given by $x \sim y$ if $x- y \in E_1$. Denote by $[x]$ the equivalent class of $x \in E$ (i.e. $[x] = x + E_1$).
\end{definition}

\begin{remark}
$E/E_1$ is a linear vector space. $x_1,x_2 \in E$ $\Rightarrow$ $[x_1] + [x_2] = [x_1 + x_2]$, $\lambda [x] = [\lambda x]$ for $x \in \E, \lambda \in \F$. i.e. $v,w \in E/E_1$, $\lambda,\mu \in \F$, then $\lambda v + \mu w \in E/E_1$. If $E/E_1$ has finite dimension, then the dimension of $E/E_1$ is called the \underline{codimension} of $E_1$ in $E$. If $E$ is finite dimension, then $\text{codim}(E_1) + \text{dim}(E_1) = \text{dim}(E)$.
\end{remark}

\vspace{3pt}
\begin{definition}[Linear operator]\ \\
A map $T: E \to F$ between two linear spaces is a \underline{linear operator} if it preserves the properties of addition and scalar multiplication (i.e. $T(\lambda v + \mu w) = \lambda T(v) + \mu T(w)$, $v,w \in E, \lambda,\mu \in \F$).
\end{definition}
\noindent The \underline{Kernal} \& \underline{Image} of $T$ are the subspaces
\begin{itemize}
    \item $\ker(T) = \{x \in E: Tx = 0\}$
    \item $\Im(T) = \{Tx \in F: x \in E\}$
\end{itemize}
$\ker(T) \subset E,\ \Im(T) \subset F$ are subspaces.

\begin{example}\ \\
$E = \{f: K \to \R: f(\cdot) \text{ continuous}\}$, $E_1 = \{f(\cdot) \in E: f(k_1) = 0\}$, where $k_1 \in F$ is fixed. Then $E/E_1 = \text{ constant functions }$, $\text{dim}(E/E_1) = 1$.
\end{example}

\clearpage
\subsection{Normed Spaces}
\begin{definition}\ \\
Let $E$ be a linear vector space, a norm $\norm{\cdot}: E \to \R,\ x \to \norm{x}$ satisfies the following properties:
\begin{itemize}
    \item $\norm{x} \geq 0,\ x \in E$ and $\norm{x} = 0 \Rightarrow x = 0$.
    \item $\norm{\lambda x} = \abs{\lambda} \norm{x},\ x \in E, \lambda \in \F$.
    \item $\norm{x + y} \leq \norm{x} + \norm{y}, x,y \in E$ (triangle inequality) 
\end{itemize}
\end{definition}
\np A linear space $E$ equipped with a norm $\norm{\cdot}$ is called a normed space. This makes $E$ a metric space with metric $d(x,y) = \norm{x -y}, x,y \in E$ (i.e. norm could generate metric).

\vspace{3pt}
\begin{definition}\ \\
A metric on E is a binary function $d(\cdot, \cdot): E \times E \to \R$ with the following properties ($x,y \in E$):
\begin{itemize}
    \item $d(x,x) = 0$ and $d(x,y) = 0$ $\Rightarrow$ $x = y$.
    \item $d(x,y) = d(y,x)$ i.e. Symmetry
    \item $d(x,z) \leq d(x,y) + d(y,z)$ (triangle inequality)
\end{itemize}
\end{definition}
\begin{examples}\ 
\begin{enumerate}[label = (\alph*)]
    \item Let $\ell_\infty$ be the space of bounded sequences $x = (x_1,x_2,\dots)$ with $x_j \in \R,\ j = 1,2,\dots$. Define $\norm{x} = \norm{x}_\infty = \sup_{j \geq 1} \abs{x_j}.$
    \item Let $\ell_1$ be the space of absolutely summarable sequences $x = (x_1,x_2,\dots)$ s.t. $\sum_{j=1}^\infty \abs{x_j} < \infty$. Define $\norm{x} = \norm{x}_1 = \sum_{j = 1}^\infty \abs{x_j} < \infty$.
    \item The space $C(K)$ or continuous functions $f: K \to \R$, where $K$ is a compact Hausdorff space. Define $\norm{f}_\infty = \sup_{x \in K} \abs{f(x)}$
\end{enumerate}
It can be seen that $\ell_1 \subset \ell_\infty$, as $\norm{x}_\infty \leq \norm{x}_1$.
\end{examples}

\vspace{6pt}
\subsection{Geometry of Normed Spaces}
\begin{definition}\ \\
A (closed) ball centered at a point $x_0 \in E$ with radius $r > 0$ is the set $B(x,r) := \{x \in E: \norm{x - x_0} \leq r\}$.\\ The sphere centered at $x_0$ with radius $r > 0$ is the set $S(x_0; r) := \{ x \in E: \norm{x - x_0} = r\}$. i.e. $S(x_0;r)$ is the "Boundary" of $B(x_0; r)$, $S(x_0;r) = \partial B(x_0;r)$. 
\end{definition}

\begin{remark}\ \\
Balls can have different geometries on the properties of the norm. Unit ball centered at origin for $\norm{\cdot}_\infty$ is just a "square" centered at $0$. Unit ball for $\norm{\cdot}_1$ is also a "square". $B(0,1) = \{x = (x_1,x_2,\dots): -1 < y_\ep < 1,\ \forall\ \ep\}$, $y_\ep = \sum_{j = 1}^\infty \ep_j x_j,\ \ep_j = \pm 1$. It can be shown that different norms give different geometry.
However, they have important common features, most notably convexity properties.
\end{remark}

\begin{definition}\ \\
A set $K \subset E$ is convex ($E$ linear space) if $\forall\ x,y \in K$, $\forall\ 0 \leq \lambda \leq 1$, $\lambda x + (1 - \lambda) y \in K$.
\end{definition}
\begin{remark}
i.e. $x,y \in E$, then the line segment joining $x$ and $y$ is in $K$.
\end{remark}

\begin{definition}\ \\
$f: E \to \R$ is a convex function if $f(\lambda x + (1 - \lambda)y) \leq \lambda f(x) + (1 - \lambda) f(y)$ for $x,y \in E,\ 0 \leq \lambda \leq 1$.
\end{definition}
\begin{remark}
If $f: E \to \R$ is a convex function, then for any $M \in \R$ the set $\{x \in E: f(x) \leq M\}$ is convex.
\end{remark}

\vspace{3pt}
\begin{proposition}\ \\
Let $(E, \norm{\cdot})$ be a normed linear space, then the function $x \to \norm{x},\ x \in E$ is convex and continuous. 
\end{proposition}
\clearpage
\begin{proof}\ \\
    Let $f: E \to \R$ be $f(x) = \norm{x}$, $f(x) - f(y) = \norm{x} - \norm{y} \leq \norm{x - y}$ by triangle inequality $\Rightarrow$ 
    \begin{equation*}
        \abs{f(x) - f(y)} \leq \norm{x - y},\ x,y \in E
    \end{equation*}
    $\Rightarrow$ $f(\cdot)$ is \underline{Lipschitz continuous}.\\
    For convexity, let $0 < \lambda < 1$, then
    \begin{align*}
        f(\lambda x + (1 - \lambda)y) &= \norm{\lambda x + (1-\lambda )y} \leq \norm{\lambda x} + \norm{(1 - \lambda)y} \\&= \lambda \norm{x} + (1 - \lambda)\norm{y}\\ &= \lambda f(x) + (1-\lambda) f(y).
    \end{align*}
\end{proof}
\begin{remark}\ \\
$f(\cdot)$ continuous implies the closed ball 
\begin{equation*}
    B(x_0; r) = \{x \in E: \norm{x - x_0} \leq r\} = \{x \in E: f(x - x_0) \leq r\}
\end{equation*} 
is closed in topology of $E$. $f$ convex $\Rightarrow$ $B(x_0;r)$ convex.
\end{remark}
\begin{remark}\ \\
$f: E \to \R$ convex $\Rightarrow$ the sub-level set $\{x \in E: f(x) \leq M\}$ is also convex. However, it is possible to have non-convex functions $f(\cdot)$ such that sets $\{x \in E: f(x) \leq M\}$ are convex.
\begin{example}\ \\
Take $f(x) = \abs{x}^p,\ x \in \R,\ p > 0$, then $f(\cdot)$ is convex if $p > 1$ and non-convex if $p < 1$. The sets $\{x \in \R: f(x) \leq M\}$ are all convex independent of $p$.
\end{example}
It is also possible to go in opposite direction if $f(\cdot)$ has the dilation property s.t. $f(\lambda x ) = \abs{\lambda} f(x)$.
\begin{corollary}\label{1.10}\ \\
Suppose $x \to \norm{x}$ satisfies
\begin{enumerate}[label = (\alph*)]
    \item $\norm{x} \geq 0, \norm{x} = 0$ $\Leftrightarrow$ $x = 0$.
    \item $\norm{\lambda x} = \abs{\lambda} \norm{x}, \lambda \in \F$.
    \item The unit ball $B(0;1)$ is convex.
\end{enumerate}
Then $F: x \to \norm{x}$ satisfies the \underline{triangle inequality}
$$\norm{x + y} \leq \norm{x} + \norm{y}.$$
\end{corollary}
\begin{proof}\ \\
    $u,v \in B(0,1),\ 0 < \lambda < 1$ $\Rightarrow$ $\lambda u + (1 - \lambda) v \in B(0,1)$. Let $x,y \in E$, $\lambda = \norm{x} / (\norm{x} + \norm{y})$ and $u = \frac{x}{\norm{x}},\ v = \frac{y}{\norm{y}}$, then
    $$\frac{\norm{x}}{\norm{x} + \norm{y}} \frac{x}{\norm{x}} + \frac{\norm{y}}{\norm{x} + \norm{y}} \frac{y}{\norm{y}} \in B(0,1)$$
$\Rightarrow$ $\norm{\frac{x}{\norm{x} + \norm{y}} + \frac{y}{\norm{x} + \norm{y}}} \leq 1$, and from (b) it follows the triangle inequality.
\end{proof}
\begin{remark}
If $x \to \norm{x}$ satisfies (a) (b) and is a convex function, then it satisfies (c).
\end{remark}
Proof of triangle inequality using convexity: $\frac{1}{2} \norm{x + y} = \norm{\frac{x}{2} + \frac{y}{2}} \leq \frac{1}{2} (\norm{x} + \norm{y})$.
\end{remark}



\vspace{6pt}
\np Let $E$ be a normed space, $E_1$ be a subspace of E, with the \underline{quotient space} $E/ E_1$. We try to define a norm on $E / E_1$ by $$\norm{[x]} = \inf_{y \in E_1} \norm{x + y}.$$
In this case, if $x \in \overline{E_1} - E_1$, then we may have $\norm{[x]} = 0$ but $[x] \neq 0 \in E/E_1$. Notice this is different from the finite dimension case. All finite dimension spaces $E_1$ are closed but not in general if $E_1$ has $\infty$-dimensions. 
\begin{example}\ \\
$\ell_1 (\R) = \{x = (x_1,x_2,\dots): x_i \in \R\}$ s.t. $\norm{x}_1 < \infty$, $\norm{x}_1 = \sum_{j = 1}^\infty \abs{x_j}$.\\ 
$E_1 = \{x = (x_1,x_2,\dots):\ x_j \in \R,\ \text{card}(\{x_j \neq 0,\ j = 1,2,\dots\}) < \infty \}$.\\ 
It can be checked that $\overline{E_1} = l_1(\R)$.
\end{example}

\begin{proposition}\ \\
Let $(E, \norm{\cdot})$ be a normed space and $E_1 \subset E$, where $E_1$ is a closed subspace. Then $\norm{\cdot}: E/E_1 \to \R$ defined by $\norm{[x]} = \inf_{y \in E_1} \norm{x+y}$ is a norm on $E/E_1$.
\end{proposition}
\begin{proof}\
    \begin{enumerate}[label = (\alph*)]
        \item $\norm{[x]} = 0$ $\Rightarrow$ $\inf_{y \in E_1} \norm{x + y} = 0$ $\Rightarrow$ $x \in E_1$ since $E_1$ closed $\Rightarrow$ $[x] = 0$.
        \item $\norm{\lambda [x]} = \inf_{y \in E_1} \norm{\lambda x + y} = \inf_{z \in E_1} \norm{\lambda x + \lambda z} = \abs{\lambda} \inf_{z \in E_1} \norm{x +z} = \abs{\lambda} \norm{[x]}$.
        \item $\norm{[x] + [y]} = \inf_{x_1,y_1 \in E_1} \norm{x + y + x_1 + y_1} \leq \inf_{x_1 \in E_1} \norm{x + x_1} + \inf_{y_1 \in E_1} \norm{y + y_1} = \norm{[x]} + \norm{[y]}$.
    \end{enumerate}
\end{proof}

\begin{remark}\ \\
In fact, since $E_1$ is a subspace, we have $[0] = E_1$, and thus $\norm{[x]} = \inf_{y \in E_1} \norm{x + y} = \inf_{x \in [x]} \norm{x}$. This is because $[x] = [x - y],\ \forall\ y \in E_1$.
\end{remark}

\clearpage
\section{Banach Spaces}
\subsection{Banach Space}
\begin{definition}[Banach Space]\ \\
A linear normed space is a \underline{Banach space} if it is complete (i.e. every Cauchy sequence converges).
\end{definition}
\begin{remark}\ \\
Completeness means if $x_n \in E,\ n \geq 1$ is a Cauchy sequence has property:
$$\lim_{m \to \infty} \sup_{n \geq m} \norm{x_n - x_m} = 0$$
then $\exists\ x_\infty \in E$ s.t. $\lim_{n \to \infty} \norm{x_n - x_\infty} = 0$.
\end{remark}
\begin{example}
The spaces $\ell_1, \ell_\infty$ and $C(K)$ are \underline{Banach spaces}.
\end{example}

\begin{remark}\ \\
Let $E$ be a linear normed space and $\{x_k\}_{k=1}^\infty$ a sequence in $E$. We say the sequence is \underline{absolutely summable} if $\sum_{k  =1}^\infty \abs{x_k} < \infty$.
\end{remark}


\vspace{3pt}
\begin{theorem}[Criterion for completeness]\ \\
A normed space $(E, \norm{\cdot})$ is a Banach space \underline{if.f} every \underline{absolutely summable} sequence in $E$ converges.
\end{theorem}
\begin{proof}\
\begin{itemize}
    \item $\Rightarrow:$\ \\ 
    Suppose $E$ is a \underline{Banach space} and $\{x_k\}_{k=1}^\infty$ a \underline{absolutely summable} sequence. Set $S_n = \sum_{k = 1}^n x_k,\ n \geq 1$. Show $S_n$ is Cauchy, and thus $E$ complete implies $\exists\ S_\infty$ s.t. $\lim_{n \to \infty} \norm{S_n - S_\infty} = 0$.\\ For $n\geq m$, 
    \begin{equation*}
        \norm{S_n - S_m} = \norm{\sum_{k = m+1}^n x_k} \leq \sum_{k = m+1}^n \norm{x_k} \leq \sum_{k = m+1}^\infty \norm{x_k}.
    \end{equation*} 
    Therefore, $\lim_{m \to \infty} \sum_{k = m+1}^\infty \norm{x_k} = 0$, which shows the sequence $\{x_k\}_{k = 1}^\infty$ is Cauchy.
    \item $\Leftarrow:$\ \\ 
    Suppose $E$ is not complete, then $\exists$ a Cauchy sequence $\{x_n\}_{n=1}^\infty$ which does not converge. Furthermore, no subsequence of $\{x_n\}_{n=1}^\infty$ converges. Construct an \underline{absolutely summable} sequence which does not converge.\\
    Define $n(1) \geq 1$ s.t. $\norm{x_n - x_{n(1)}} \leq \frac{1}{2}$ $\forall\ n > n(1)$, and let $n(2) > n(1)$ be s.t. $\norm{x_n - x_{n(2)}} \leq \frac{1}{2^2}$ $\forall\ n > n(2)$. Construct let this and get $n(1), n(2), n(3),\dots$ s.t. $\norm{x_n - x_{n(k)}} \leq \frac{1}{2^k},\ \forall\ n > n(k)$. Define $w_j = x_{n(j+1)} - x_{n(j)},\ j = 1,2,\dots$, then
    $$x_{n(m)} = x_{n(1)} + \sum_{j = 1}^{m-1} w_j,\ m = 1,2,\dots$$
    $\{x_{n(m)}\}$ does not converge. Hence, series $\sum_{j = 1}^{m-1} w_j$ does not converge. However,
    $$\sum_{j = 1}^\infty \norm{w_j} \leq \sum_{j = 1}^\infty \frac{1}{2^j} = 1$$
    which shows $\{w_j\}$ is \underline{absolutely summable}, and thus $\sum_{j = 1}^m w_j$ converges, which reachs a contradiction.
\end{itemize} 
\end{proof}

\clearpage
\subsection{Completion of normed spaces to a Banach space}
\begin{theorem}\ \\
Suppose $E$ is a normed space, then $\exists$ a \underline{Banach space} $\hat{E}$ called a completion of $E$ with the following properties:\\
$\exists$ a linear map $i: E \to \hat{E}$ s.t. 
\begin{enumerate}[label = (\alph*)]
    \item $\norm{ix} = \norm{x},\ x \in E$.
    \item $\Im(i)$ is dense in $\hat{E}$.
\end{enumerate}
\end{theorem}
\begin{proof}
The core of the proof has been presented in HW1 Pr4, and the remaining thing is to show $\hat{E}$ is a \underline{Banach space}. \textcolor{red}{Put them here later}.

\np Denote $(x_n)$ the Cauchy sequence in $E$, and $[(x_n)^N]$ the Cauchy sequence in $\hat{E}$. For each $N$, $\exists\ n_N > 0$ s.t. $\norm{x_n^N - x_m^N}_E \leq \frac{1}{N},\ \forall\ m,n \geq n_N$. Choose $y_N := x_{n_N}^N$, we construct the sequence $[y_N] = [(x_{n_N})^N]$. Note that $\forall\ \ep > 0$, $\exists\ N_\ep$ s.t.
\begin{equation*}
    \norm{(x_n)^N - (x_n)^M}_{\hat{E}} = \lim_{n \to \infty} \norm{x_n^N - x_n^M}_E  \leq \ep,\ \forall\ N,M \geq N_\ep \geq \frac{1}{\ep}.
\end{equation*}
Thus we can further find $n_\ep > 0$ such that
\begin{equation*}
    \norm{x_n^N - x_n^M}_E \leq 2 \ep,\ \forall\ n \geq n_\ep.
\end{equation*}
Taking $n = \{n_\ep,n_N,n_M\}$, we have
\begin{equation*}
    \norm{y_N - y_M}_E \leq \norm{y_N - x_n^N}_E + \norm{x_n^N - x_n^M}_E + \norm{x_n^M - y_M}_E \leq\frac{1}{N} + 2\ep  + \frac{1}{M} \leq  4\ep.
\end{equation*}
therefore, $\forall\ \ep > 0$, $\exists\ N_\ep$ s.t. $\norm{y_N - y_M}_E \leq 4\ep,\ \forall\ N,M \geq N_0$, which shows $[y_N]$ is a Cauchy sequence in $E$. Finally we want to show that $[(x_n)^N]$ converges to $[y_N]$. Given $\ep > 0$, can choose $n_\ep > 0$ s.t.
\begin{equation*}
    \norm{y_n - y_m} < \ep,\ \forall\ n,m \geq n_\ep,
\end{equation*}
then we can pick $N_0 \geq \max{\{1/\ep, n_\ep\}}$. For all $N \geq N_0$,
\begin{equation*}
    \norm{x_n^N - y_N}_E = \norm{x_n^N - x_{n_N}^N}_E \leq \frac{1}{N} \leq \frac{1}{N_0} \leq \ep,\ \forall\ n \geq n_N. 
\end{equation*}
Fix $N \geq N_0$, then
\begin{equation*}
    \norm{x_n^N - y_n}_E \leq \norm{x_n^N - y_N}_E + \norm{y_N - y_n}_E \leq 2\ep,\ \forall\ n \geq \max{\{n_N,n_\ep\}}.
\end{equation*}
Therefore,
\begin{equation*}
    \norm{[(x_n)^N] - [y_n]}_{\hat{E}} = \lim_{n \to \infty} \norm{x_n^N - y_n}_E \leq 2\ep
\end{equation*}
Since $\ep > 0$ is arbitrary, we complete the proof.
\end{proof}


\begin{remark}\ \\
    $i$ is an \underline{isometric} embedding of $E$ into $\hat{E}$ (Property (a)). $\hat{E}$ is the smallest \underline{Banach space} containing image of $i$ over $E$ (Property (b)). 
\end{remark}


\vspace{3pt}
\begin{theorem}[$\ell_p$ space is complete]\ \\
For $1 \leq p \leq \infty$, $\ell_p$ space is a \underline{Banach space}.
\end{theorem}
\begin{proof}\ \\
Suppose $\{x^n\}_{n \in \N}$ is a Cauchy sequence in $\ell_p$, where $x^n = \{x^n_k\}_{k \in \N}$. Then $\forall\ \vep > 0$, $\exists\ N \in \N$ s.t. (for $1 \leq p < \infty$, and $p = \infty$ case is trivial)
\begin{equation*}
    \norm{x^n - x^m}_p = (\summ{k=1}{\infty} \abs{x^n_k - x^m_k}^{p})^{1/p} < \vep,\ \forall\ n > m \geq N.
\end{equation*}
The above equation implies that for fixed $k$, $\{x_k^n\}_{n \in \N}$ forms a Cauchy sequence in $\R$, and thus $x_k^n \to x_k^\infty$ for some $x^\infty_k \in \R$, which implies $x^n \to x^\infty$. It remains to show
\begin{itemize}
    \item $x^\infty \in \ell_p:$\\
    Every Cauchy sequence is bounded, and thus $\exists\ M > 0$ s.t.
    \begin{equation*}
        \norm{x^n}_p = (\summ{k=1}{\infty} \abs{x^n_k}^p)^{1/p} \leq M,\ \forall\ n \in \N.
    \end{equation*}
    Therefore, we have
    \begin{equation*}
        \summ{k=1}{m} \abs{x^n_k}^p \leq M^p,\ \forall\ m,n \in \N.
    \end{equation*}
    Take $n \to \infty$ and then $m \to \infty$, we complete the claim.
    \item $x^n \to x^\infty$ in $\ell_p$:\\
    Fix $n > N$, and note that $\summ{k=1}{\infty} \abs{x^n_k - x^m_k} < \vep^p,\ \forall\ m > n$, which implies for each $K$,
    \begin{equation*}
        \summ{k=1}{K} \abs{x^n_k - x^m_k} < \vep^p.
    \end{equation*}
Take $m \to \infty$, we have
    \begin{equation*}
        \summ{k=1}{K} \abs{x^n_k - x^\infty_k} < \vep^p.
    \end{equation*}
    Finally take $K \to \infty$, we complete the proof.
\end{itemize}
\end{proof}
\begin{remark}\ \\
    For $1 < p < \infty$, $x = \{x_n\}_{n = 1}^\infty \in \ell_p$ if $\sum_{n=1}^\infty \abs{x_n}^p < \infty$, $\norm{x}_p = (\sum_{n=1}^\infty \abs{x_n}^p)^{1/p}$.\\
    Show that $x \to \norm{x}_p$ satisfies properties of a norm:
    \begin{itemize}
        \item $\norm{x} \geq 0$, $\norm{x} = 0$ $\Leftrightarrow$ $x = 0$ (easy).
        \item $\norm{\lambda x} = \abs{\lambda} \norm{x}, \lambda \in \F$ (easy).
        \item Triangle inequaility: Recall \hyperref[1.10]{Cor 1.10}, it suffices to show $B(0,1)$ is convex. Note $f(x) = \sum_{n=1}^\infty \abs{x_n}^p$ is a convex function, since $x \to \abs{x}^p,\ x \in \R$ is convex for $p \geq 1$.
    \end{itemize}
    Hence $\norm{x + y}_p \leq \norm{x}_p + \norm{y}_p$ and
        \begin{equation*}
            (\sum_{j=1}^\infty \abs{x_j + y_j}^p)^{1/p} \leq (\sum_{j=1}^\infty \abs{x_j}^p)^{1/p} + (\sum_{j=1}^\infty \abs{y_j}^p)^{1/p} \quad (\text{Minkowski Inequality})
        \end{equation*}
    The usual proof of \underline{Minkowski Inequality} is via the \underline{H\"older Inequality}.
    \begin{lemma}[H\"older Inequality for $\ell_p$ space]
    \begin{equation*}
        \sum_{j = 1}^\infty \abs{x_j y_j} \leq (\sum_{j = 1}^\infty \abs{x_j}^p)^{1/p} (\sum_{j = 1}^\infty \abs{y_j}^q)^{1/q}
    \end{equation*}
, where $\frac{1}{p} + \frac{1}{q} = 1.$
    \end{lemma}
\begin{proof}
\begin{itemize}
    \item Consider $F(t) = t - \frac{t^p}{p},\ t \geq 0$, we have by finding the maximum
    \begin{equation*}
        t \leq \frac{t^p}{p} + 1 - \frac{1}{p} = \frac{t^p}{p} + \frac{1}{q},\ \frac{1}{p} + \frac{1}{q} = 1.
    \end{equation*}
    \item Taking $t = \frac{\alpha}{\beta^{q - 1}}$, we have
    \begin{equation*}
        \alpha \beta \leq \frac{\alpha^p}{p} + \frac{\beta^q}{q}.
    \end{equation*}
    which is known as \underline{Yound's Inequality}.
    \item WLOG, assume $0 < \norm{f}_p, \norm{g}_q < \infty$. Consider $F(x) = \frac{f(x)}{\norm{f}_p},\ G(x) = \frac{g(x)}{\norm{g}_p}$. Note $\norm{F}_p = \norm{G}_p = 1$, we have
    \begin{align*}
        \int \abs{F(x) G(x)} d\mu &\leq \int \frac{\abs{F(x)}^p}{p} + \frac{\abs{G(x)}^q}{q} d\mu\\
        \frac{\norm{fg}_1}{\norm{f}_p \norm{g}_q} &\leq \frac{1}{p} + \frac{1}{q} = 1
    \end{align*}
    which is a more general case of \underline{H\"older Inequality}. 
\end{itemize}
\end{proof}
\begin{remark}\ \\
    The \underline{Young's Inequality} is a special case of the inequality $$xy \leq f(x) + \Ls f(y)$$
where $\Ls f(y) := \sup_x [xy - f(x)]$ is the \underline{Legendre Transform} of $f(\cdot)$.\\
If $f(\cdot)$ is convex, then $F_N: (x,y) \to xy - f(x)$ is concave and has unique maximum.
\end{remark}
\begin{remark}\ \\
    $\Ls f(\cdot)$ is always convex even if $f(\cdot)$ is not. $f(x) = \frac{x^p}{p}$ then $\Ls f(y) = \frac{y^{p'}}{p'}$
\end{remark}
\begin{proof}[Proof of \underline{Minkowski Inequality}]\
\begin{equation*}
    \sum_{j = 1}^\infty \abs{x_j + y_j}^p \leq \sum_{j=1}^\infty \abs{x_j} \abs{x_j + y_j}^{p-1} + \sum_{j= 1}^\infty \abs{y_j} \abs{x_j + y_j}^{p-1}
\end{equation*}
and by H\"older inequality, we have 
\begin{equation*}
    \sum_{j = 1}^\infty \abs{x_j} \abs{x_j + y_j}^{p-1} \leq (\sum_{j= 1}^\infty  \abs{x_j}^p)^{1/p} (\sum_{j = 1}^\infty \abs{x_j + y_j}^{(p - 1) p'})^{1 /p'},
\end{equation*}
where $\frac{1}{p} + \frac{1}{p'} = 1$. Then proceed to complete the proof.
\end{proof}
\end{remark}

\np Let $(\Omega, \Sigma,\mu)$ be a measure space and $\Ls_p(\Omega, \Sigma, \mu)$ be all measurable functions $f: \Omega \to \F$ s.t. $\int_\Omega \abs{f}^p \dr\mu < \infty$. Then $\Ls_p (\Omega, \Sigma, \mu)$ is a normed space with norm $\norm{f}_p = (\int_\Omega \abs{f}^p \dr \mu)^{1/p}$.

\vspace{3pt}
\begin{theorem}\ \\
$\Ls_p(\Omega, \Sigma, \mu)$ is a \underline{Banach space}.
\end{theorem}
\begin{proof}\ \\
Let $\{f_n\}_{n = 1}^\infty$ be an \underline{absolutely summable} sequence in $\Ls^p$, then $\norm{\sum_{k = 1}^N f_k}_p \underbrace{\leq}_{\text{Minkowski}} \sum_{k = 1}^N \norm{f_k}_p \leq C $ (Independent of $N$). Hence $\int_\Omega \abs{\sum_{k=1}^N f_k}^p \dr\mu \leq C^p$.\\
\textbf{Case 1:} Assume all $f_k$ is non-negative. then the Monotone Convergence Theorem implies
$$\lim_{n \to \infty} \int_\Omega (\sum_{k=1}^N f_k)^p \dr\mu = \int_\Omega (\sum_{k=1}^\infty f_k)^p \dr\mu \leq C^p$$
Hence $g = \sum_{k=1}^\infty f_k \in \Ls_p$. Then we need to show $\sum_{k=1}^N f_k \to g$ in $\Ls_p$. Set $S_n = \sum_{k = n + 1}^\infty f_k$, then $S_n$ is a decreasing sequence $S_n \to 0$ a.e.. Conclude $\lim_{n \to \infty} \norm{S_n}_p = 0$ by DCT, as $\int_\Omega S_1^p \dr \mu < \infty$.\\
\textbf{Case 2:} For arbitrary $f_k: \Omega \to \R$, write $f_k = f_k^+ + f_k^-,\ f_k^+ = \sup{(f_k, 0)},\ f_k^- = \inf(f_k, 0)$. The sequence $\{f_k^+\}_{k = 1}^\infty$ is \underline{absolutely summable}. Proceed as before similarly.
\end{proof}

\clearpage
\section{Hilbert Spaces}
\subsection{Inner Product Space}
\begin{definition}\ \\
Let $E$ be a linear space over $\C$. An inner product 
\begin{equation*}
\inprod{\cdot}{\cdot}: E \times E \to \C     
\end{equation*}
is a function which has the following properties:
\begin{enumerate}[label = (\alph*)]
    \item $\inprod{x}{x} \geq 0,\ x \in E$ and $\inprod{x}{x} = 0$ \underline{if.f} $x = 0$.
    \item $\inprod{ax + by}{z} = a\inprod{x}{z} + b \inprod{y}{z}$.
    \item $\inprod{x}{y} = \overline{\inprod{y}{x}}, x,y \in E$.
\end{enumerate}
\end{definition}
\begin{remark}
    Inner products of spaces over $\R$ is defined similarly, but no complex conjugation in $(c)$.
\end{remark}

\vspace{3pt}
\begin{definition}[Orthogonality]\ \\
$x,y \in E$ is orthogonal if $\inprod{x}{y} = 0$, denoted by $x \perp y$.
\end{definition}

\vspace{3pt}
\begin{theorem}[Cauchy-Schwarz Inequality]\ \\
$x,y \in E$, then $\abs{\inprod{x}{y}} \leq \inprod{x}{x}^{\frac{1}{2}} \inprod{y}{y}^{\frac{1}{2}}$.
\end{theorem}
\begin{proof}\ \\
Define $Q(t)$ by $Q(t) = \inprod{x + ty}{x + ty} = \inprod{y}{y} t^2 + 2t \Re \inprod{x}{y} + \inprod{x}{x}$ if $t \in \R$. $Q(t) \geq 0$ if $t \in \R$. If $Q(t) = 0$, no more than one real root $\Rightarrow$ $[Re\inprod{x}{y}]^2 \leq \inprod{x}{x} \inprod{y}{y}$. The \underline{C-S Inequality} follows by choosing $\theta \in \R$ s.t. $\inprod{x}{y} = Re\inprod{x e^{i \theta}}{y}$. (as $\inprod{x e^{i \theta}}{x e^{i \theta}} = (x e^{i\theta})^\trs (\bar{x e^{i\theta}}) = \inprod{x}{x}$)
\end{proof}

\vspace{3pt}
\begin{corollary}\ \\
The function $x \to \norm{x} = \inprod{x}{x}^{\frac{1}{2}}$ is a norm on $E$.
\end{corollary}
\begin{proof}\ \\
Triangle inequality is a consequence of the \underline{C-S Inequality}. 
\begin{align*}
    \norm{x + y}^2 &= \inprod{x+y}{x+y}\\ &= \norm{x}^2 + 2 Re\inprod{x}{y} + \norm{y}^2\\ &\underbrace{\leq}_{\text{C-S}} \norm{x}^2 + 2 \norm{x} \norm{y} + \norm{y}^2\\ &= [\norm{x} + \norm{y}]^2
\end{align*}
Other conditions are easy to check.
\end{proof}

\begin{examples}\ 
\begin{itemize}
    \item The space $\ell_2$ of square summable sequences. $x = \{x_i\}_{i \in \N},\ y = \{y_i\}_{i \in \N}$, and $\inprod{x}{y} = \sum_{j = 1}^\infty x_j \bar{y}_j$.
    \item Similarly with $\Ls_2(\Omega,\Sigma,\mu)$, $f,g \in \Ls_2(\Omega, \Sigma, \mu),\ \inprod{f}{g} = \int_\Omega f(x )\bar{g}(x) \dr\mu(x)$.
    \item The space of $m \times n$ matrices $A = (a_{i,j}),\ 1 \leq i \leq m,\ 1 \leq j \leq n$. $\inprod{A}{B} = \trace{(A B^*)}$, where $B^* = \text{Hermitian adjoint of } B$. $B = \{b_{i,j}\}$ then $B^* = \{b_{i,j}^*\}$ s.t. $b_{i,j}^* = \overline{b}_{j,i}$. The norm corresponds to this inner product is $\norm{A}_{HS} = (\sum_{ij = 1}^\infty \abs{a_{ij}}^2)^{\frac{1}{2}}$, known as the \underline{Hilbert-Schmidt norm}.
\end{itemize}
\end{examples}

\np For an inner product space, the inner product can be expressed in terms of the norm.


\begin{proposition}\ \\
Let $E$ be an inner product space, then 
\begin{enumerate}[label = (\alph*)]
    \item $\norm{x + y}^2 + \norm{x - y}^2 = 2\norm{x}^2 + 2 \norm{y}^2$ (Parallelogram Law)
    \item $\inprod{x}{y} = \frac{1}{4} \{\norm{x + y}^2 - \norm{x - y}^2 + i \norm{x + iy}^2 - i \norm{x - iy}^2\}$ (Polarization Inequality)
\end{enumerate}
\end{proposition}
\begin{proof}\
\begin{enumerate}[label = (\alph*)]
    \item \begin{align*}
        \norm{x+y}^2 + \norm{x - y}^2 &= \inprod{x + y}{ x + y} + \inprod{x - y}{x - y}\\
        &= \inprod{x}{x} + \inprod{y}{y} +  \inprod{x}{y} + \inprod{y}{x} + \inprod{x}{x} + \inprod{y}{y} - \inprod{x}{y} - \inprod{y}{x}\\
        &= 2 \norm{x}^2 + 2 \norm{y}^2
    \end{align*}
    \item \begin{align*}
        RHS &= \frac{1}{4} \{\norm{x + y}^2 - \norm{x - y}^2 + i \norm{x + iy}^2 - i \norm{x - iy}^2\}\\
        &= \frac{1}{4} [\inprod{x+y}{x+y} - \inprod{x - y}{x - y} + i \inprod{x + iy}{x + iy} - i \inprod{x - iy}{x - iy}]\\
        &= \frac{1}{4}[2 \inprod{x}{y}  + 2 \inprod{y}{x} + 2i\inprod{x}{iy} + 2 i \inprod{iy}{x}]\\
        &= \inprod{x}{y} = LHS
    \end{align*}
\end{enumerate}
\end{proof}
\begin{remark}
    Polarization identity shows that the function $x \to \norm{x}^2$ determines the inner product.
\end{remark}


\vspace{12pt}
\subsection{Hilbert Space}
\begin{definition}\ \\
A \underline{complete inner product space} is called a \underline{Hilbert space}. 
\end{definition}
\begin{remark}
    Have already seen that $\ell_2$ and $\Ls_2 (\Omega,\Sigma,\mu)$ are complete. They are \underline{Hilbert spaces}.
\end{remark}
\begin{remark}
    The key notion in \underline{Hilbert space} is \underline{orthogonality}.
\end{remark}

\vspace{3pt}
\begin{definition}\ \\
Let $A \subset \Hs := \text{ \underline{Hilbert space }}$. The \underline{orthogonal complement} $A^{\perp}$ of $A$ is 
\begin{equation*}
    A^\perp := \{x \in \Hs: \inprod{x}{y} = 0,\ \forall\ y \in A\}.
\end{equation*}
\end{definition}
\begin{remark}\ \\
    If $A^\perp$ is a closed linear subspace of $\Hs$, then $A^\perp$ is also a \underline{Hilbert space}, where closure follows from continuity of the function $x \to \inprod{x}{y},\ x \in \Hs$.
\end{remark}

\vspace{3pt}
\begin{theorem}[Orthogonality Principle]\ \\
Assume $E \subset \Hs$ is a closed linear subspace of the \underline{Hilbert space} $\Hs$ and $x \in \Hs$. Then
\begin{enumerate}[label = (\alph*)]
    \item $\exists$ a unique closed point $y = P_E x \in E$ to x. (i.e. $\norm{x - P_E x} = \inf{\norm{x - y'}},\ y' \in E$)
    \item The point $y = P_E x \in E $ is the unique vector s.t. $x - y \in E^\perp$.
\end{enumerate}
\end{theorem}
\begin{proof}\ \\
Note the function $y' \to \norm{x - y'},\ y' \in E$ is convex. We expect a minimizing $y'$, to show whose existence we typically need:
\begin{itemize}
    \item Compactness properties.
    \item Non-degeneracy Properties for Uniqueness.
\end{itemize}
Here we will use \underline{Parallelogram Law} so \underline{Compactness} is no longer needed.
\begin{itemize}
    \item Let $(y_n) \in E,\ n=1,2,\dots$ be a minimizing sequence (i.e. $\lim_{n \to \infty} \norm{x - y_n} = \inf_{y' \in E} \norm{x - y'} = d$). Note
\begin{equation*}
    \norm{y_n - y_m}^2 + 4 \norm{x - \frac{1}{2} (y_n + y_m)}^2 = 2 \norm{x - y_n}^2 + 2 \norm{x - y_m}^2
\end{equation*}
As $n,m \to \infty$, RHS $\to 4d^2$. \\
$\frac{1}{2} (y_n + y_m) \in E$ $\Rightarrow$ $\norm{x - \frac{1}{2} (y_n + y_m)} \geq d$. Conclude $\lim_{m\to \infty} \sup_{m \geq n} \norm{y_n - y_m}^2 = 0$ $\Rightarrow$ $(y_n)$ is a Cauchy sequence. Since $\Hs$ is complete, we have $y_n \to y_\infty \in E,\ \norm{x - y_\infty} = d$ by the fact that $E$ is closed. Set $P_E x = y_\infty$. $y_\infty$ is unique: if $\norm{x - y_\infty} = \norm{x - y_\infty'} = 0$, by \underline{Parallelogram Law}, $\norm{y_\infty - y_\infty'} = 0$. Hence $P_E x$ is uniquely defined.
    \item Show $P_E x$ is the unique vector $y \in E$ s.t. $x - y \in E^\perp$. Let $y' \in E$ and $Q(t)$ be the quadratic
    \begin{align*}
        Q(t) &= \inprod{x - P_E x + ty'}{x - P_E x + t y'} \\&= \norm{x - P_E x + ty'}^2
    \end{align*}
    $t \to Q(t)$ has a \underline{strict} minimum at $t = 0$\\ 
    $\Rightarrow$ $Q'(0) = 0,\ \Re{\inprod{x - P_E x}{y'}} = 0,\ \forall\ y' \in E$ $\Rightarrow$ $\inprod{x - P_E x}{y'} = 0,\ \forall\ y' \in E$.\\
    $\Rightarrow$ $x - P_E x \in E^\perp$.\\
    Finally $P_E x \in E$ is the unique vector s.t. $x - P_E x \in E^\perp$ because for such $P_E x$, we have 
    \begin{equation*}
        Q(t) = \norm{x - P_E x}^2 + t^2 \norm{y'}^2,\ \forall\ y' \in E.
    \end{equation*}
    and thus the solution is unique.
\end{itemize}
\end{proof}
\begin{remark}\ \\
    The minimizer for the function $y' \to \norm{x - y'},\ y' \in E$ is characterized by the \underline{orthogonality condition} $x - y \perp E$ for some $y \in E$.
\end{remark}

\vspace{3pt}
\begin{definition}[Orthogonal Projection]\ \\
Let $\Hs$ be a \underline{Hilbert space}, $E \subset \Hs$ a closed subspace. The projection $P_E: \Hs \to E$ is given by $x \to P_E x$, where $P_E x$ is defined uniquely to be $x - P_E x \in E^\perp$. 
\end{definition}

\vspace{3pt}
\begin{definition}\ \\
A mapping $A: \Bs \to \Bs$ on a \underline{Banach space} $\Bs$ is linear if $A(ax + by) = a\ Ax + b\ Ay,\ x,y \in \Bs,\ a,b \in \C$. The operator $A$ is \underline{bounded} if $\norm{A} = \sup_{\norm{x} = 1} \norm{Ax} < \infty$.
\end{definition}
\begin{remark}\ \\
    Note that $\norm{Ax} \leq \norm{A} \norm{x},\ x \in \Bs$. $P_E: \Hs \to E$  is a \underline{bounded linear operator} with the properties 
    \begin{equation*}
        P^2_E = P_E,
    \end{equation*}
    and $\norm{x}^2 = \norm{P_E x}^2 + \norm{(I - P_E) x}^2$ since $(I - P_E)x \perp P_E x$ $\Rightarrow$ $\norm{P_E} \leq 1,\ \norm{I - P_E} \leq 1$. In fact, $\norm{P_E} = \norm{I - P_E} = 1$, as $I - P_E$ is the orthogonal projection onto $E^\perp$.
\end{remark}

\vspace{12pt}
\subsection{Orthogonal Systems}

\begin{definition}\ \\
For a sequence $\{x_k\}_{k \geq 1}$ of non-zero vectors in \underline{Hilbert Space} $\Hs$. We say it is \underline{orthogonal} if $\inprod{x_k}{x_l} = 0$ for $k \neq l$. If in addition $\norm{x_k} = 1,\ l=1,2,\dots$, then the system is \underline{orthonormal}.
\end{definition}

\begin{examples}\
\begin{enumerate}[label = (\alph*)]
    \item $x_k = (0,0,\dots,\underbrace{1}_{\text{kth}},0,\dots,0) \in \ell_2,\ k = 1,2,\dots$ is O-N in $\ell_2$.
    \item For $L_2([-\pi,\pi])$, $\ell_k (t) = \frac{1}{\sqrt{2\pi}} e^{ikt},\ k \in \E$ is O-N. It is the Fourier basis associated with Fourier series 
    $$f(t) = \sum_{k =-\infty}^\infty a_k \frac{1}{\sqrt{2\pi}} e^{ikt},$$
    where $a_k = \frac{1}{\sqrt{2\pi}} \int_{-\pi}^\pi f(t) e^{-ikt} \dr t$.
    \begin{remark}\ \\
    Can generalize Fourier series to any Hilbert space $\Hs$. Let $\{x_k\}_{k \geq 1}$ be an O-N set in $\Hs$. For $n = 1,2,\dots$, can define 
    $$S_n(x) = \sum_{k = 1}^n \inprod{x}{x_k} x_k,\ x \in \Hs,\ S_n: \Hs \to E_n,$$
    where $E_n = \text{span}\{x_1,\dots,x_n\}$. $S_n$ is a linear operator $S_n = P_{E_n} = \underline{\text{ orthogonal projection onto } E_n}$ since $\inprod{x - S_n(x)}{x_k} = 0,\ k = 1,\dots,n$. $S_n(x) \in E_n,\ x - S_n(x) \perp E_n,\ \forall\ x \in \Hs$. Note that
    $$\norm{S_n x}^2 = \sum_{k=1}^n \abs{\inprod{x}{x_k}}^2,\ S_n = P_{E_n},$$ 
    we have $\norm{P_{E_n}x}^2 \leq \norm{x}^2$ $\Rightarrow$ $\sum_{k = 1}^\infty \abs{\inprod{x}{x_k}}^2 \leq \norm{x}^2$ for $x \in \Hs$, which is called \textit{Bessel's Inequality}.
    \end{remark}
\end{enumerate}
\end{examples}

\vspace{3pt}
\begin{theorem}\ \\
Let $\{x_k: k \geq 1\}$ be an O-N sequence in $\Hs$. Then the corresponding Fourier expansion 
\begin{equation*}
S_n(x) = \sum_{k = 1}^n \inprod{x}{x_k} x_k    
\end{equation*} 
converges, i.e. $\lim_{n \to \infty} S_n x = S_\infty x$ exists for $x \in \Hs$. 
\end{theorem}
\begin{proof}\ \\
It suffices to show sequence $S_n x,\ n = 1,2,\dots$ is Cauchy for each $x \in \Hs$. We have 
$$\norm{S_n x - S_m x}^2 = \sum_{k = m+1}^n \abs{\inprod{x}{x_k}}^2.$$ 
By Bessel's Inequality, $\forall\ \ep > 0$, 
\begin{equation*}
    \exists\ m(\ep) \st \sum_{k = m(\ep) + 1}^\infty \abs{\inprod{x}{x_k}}^2 \leq \ep \Rightarrow \norm{S_n x - S_{m(\ep)} x}^2 < \ep \text{ if } n > m(\ep)
\end{equation*} 
Hence $S_n x,\ n \geq 1$ is Cauchy \imply $\lim_{n \to \infty} S_n x = S_\infty x \in \Hs$. $S_\infty = P_{E_\infty}$, where $$E_\infty = \text{ closure of the linear space generated by the sequence } \{x_k\}_{k \geq 1}.$$
\begin{remark}\ \\
    $S_n = P_{E_n}$, $E_n = $ space spanned by $\{x_1,\dots,x_n\}$. $\exists\ S_\infty: \Hs \to \Hs$, $\lim_{n \to \infty} \norm{S_\infty (x) - S_n(x)} = 0,\ x \in \Hs$, and $S_\infty = P_{E_\infty}$ where $E_\infty = $ closure of the space spanned by $\{x_1,\dots,x_n\}$. Here we take closure because there are non-closed subspace of a Hilbert space.
\end{remark}
\begin{example}\ \\
Take $V$ to be the space of sequences $(a_1,a_2,\dots)$ of real numbers with norm $\norm{(a_1,a_2,\dots)} = \summ{k=1}{\infty} \abs{a_k}$. Consider the subspace of $W \subset V $ consisting of sequences with only finitely many of $a_1,a_2,\dots$ being non-zero. Then $A_k=(1,\frac{1}{2},\dots,\frac{1}{2^k},0,0,\dots), k=1,2,\dots,$ gives a sequence of elements of W, but the limit of this sequence is not in W.
\end{example}
\end{proof}


\vspace{3pt}
\begin{definition}
A system of vectors $\{x_k\}_{k \geq 1}$ in \underline{Hilbert space} $\Hs$ is \underline{complete} if the space spanned by $\{x_k\}_{k \geq 1}$ is \underline{dense} in $\Hs$.
\end{definition}
\begin{remark}\ \\
    If an \underline{O-N} set $\{x_k\}_{k \geq 1}$ is complete then we have $E_\infty = \Hs,\ P_{E_\infty} = I$ $\Rightarrow$ $x = \sum_{k=1}^\infty \inprod{x}{x_k}x_k,\ x \in \Hs$. This is \underline{Fourier Inversion Formula}. 
\end{remark}
\begin{remark}\ \\
Recall $\norm{x}^2 = \norm{P_{E_n} x}^2 + \norm{(I - P_{E_n}) x}^2$. Let $n \to \infty$ $\Rightarrow$ $\norm{x}^2 = \lim_{n \to \infty} \norm{P_{E_n} x}^2 = \sum_{k=1}^\infty \abs{\inprod{x}{x_k}}^2$. $\norm{x}^2 = \sum_{k=1}^\infty \abs{\inprod{x}{x_k}}^2$ preserves identity. 
\end{remark}


\vspace{3pt}
\begin{definition}\ \\
A metric space is \underline{separable} if it contains a countable dense subset.
\end{definition}
\begin{remark}\ \\
For Banach space separability follows from finding a countable set of vectors $\{x_k\}_{k \geq 1}$ s.t. span of the $\{x_k\}_{k \geq 1}$ is dense in $E$.
\end{remark}


\vspace{12pt}
\subsection{Gram-Schmidt Orthogonalization}

\np Suppose $x_1,x_2,\dots \in \Hs$ is a set of vectors and $E_n = \text{span}\{x_1,\dots,x_n\}$. Then we can find an \underline{O-N} set $\{y_k\}_{k \geq 1}$ s.t. $E_n = \text{span}\{y_1,\dots,y_{m(n)}\}$ where $m(n) \leq n$. Set $y_1 = \frac{x_1}{\norm{x_1}}$, $y_n = \frac{(I - P_{E_{n-1}} ) x_n}{\norm{(I - P_{E_{n-1}}) x_n}}$ if $x_n \not\in E_{n-1}$, i.e. $E_{n-1}$ is properly contained in $E_n$.
\begin{remark}\ \\
Proving completeness of a set of vectors $\{x_k\}_{k \geq 1}$ in $\Hs$ can be non-trival.
\end{remark}

\begin{examples}\
\begin{enumerate}[label = (\alph*)]
    \item Haar basis for $\Ls^2([0,1])$: Let $h: (0,1) \to \R$ s.t.
    \begin{equation*}
        h(t) = 1 \text{ for } t < \frac{1}{2},\ h(t) = -1 \text{ for } \frac{1}{2} < t < 1.
    \end{equation*}
     Extend $h(\cdot)$ by zero outside $(0,1)$ $\Rightarrow$ $h: \R \to \R$, $h(t) = 0$ if $t \notin (0,1)$. The function $t \to h(2^k t)$ has support in interval $0 < t < 2^{-k}$. Move this support to $l 2^{-k} < t < (l+1)2^{-k}$. By translation, set 
     $$h_{k,l} (t) = h(2^kt - l),\ l=0,1,2,\dots, 2^k -1.$$ 
     The constant function plus functions $h_{k,l},\ k=0,1,\dots,\ 0 \leq l \leq 2^k-1$ are a complete orthogonal set for $\Hs = \Ls^2([0,1])$.
    \begin{proof}\ \\
    The span of the functions includes characteristic functions $\1_F$ for any dyadic intervals \begin{equation*}
    [2^{-k} l, 2^{-k} (l+1)],\ l=0,\dots,2^{k}-1,\ k=1,2,\dots.    
    \end{equation*} 
    If the set is \underline{not} complete, then $\exists\ f \in \Ls^2([0,1])$ s.t. $\int_F f dt = 0$ for all dyadic intervals $F$ and $f$ does not vanish a.e.. Can approximate any measurable sets $E \subset (0,1)$ by a union of dyadic intervals. An easy way to see this is $\{F \in \Bs: \int_F f dt = 0\} = \Bs((0,1))$. Then observe dyadic intervals generate all open intervals. Hence $\int_F f dt = 0$ for all measurable $F \subset (0,1)$. Let $F_+ = \{t \in (0,1): f(t) > 0\}$, if $m(F_+) > 0$, then $\int_{F_+} f dt > 0$. Hence $m(F_+) = 0$, and similarly can show $m(F_-) = 0$ \imply $f = 0$ a.e.. Complete the proof. 
    \end{proof}
    
    \item Fourier basis $e_k(t) = \frac{1}{\sqrt{2\pi}} e^{ikt},\ k \in \R,\ -\pi < k < \pi$ is complete in $\Ls^2([-\pi,\pi])$.
    
    \begin{proof}\ \\
    Use \underline{Stone-Weierstrass Thm}. Application to Fourier series. All $e_k(\cdot)$ are in $C([-\pi,\pi])$, i.e. continuous functions $f: [-\pi,\pi] \to \C$. $C([-\pi,\pi])$ is a Banach space with sup norm $\norm{f} = \sup_{t \in [-\pi,\pi]} \abs{f(t)}$. \underline{Stone-Weierstrass Thm} implies density of the $e_k(\cdot),\ k \in \R$ in $C([-\pi,\pi])$. Completeness in $\Ls^2([-\pi,\pi])$ follows from the density of continuous functions in $\Ls^2([-\pi,\pi])$. 
    \end{proof}
    \begin{remark}[Stone-Weierstrass Theorem]\ \\
    Suppose $f$ is a continuous real-valued function defined on the real interval $[a, b]$. For every $\ep > 0$, there exists a polynomial $p$ such that for all $x \in [a,b]$, we have $\abs{f(x) - p(x)} < \ep$, or equivalently, the supremum norm $\norm{f - p}_\infty$ < $\ep$.
    \end{remark}
\end{enumerate}
\end{examples}















\clearpage
\section{Bounded Linear Functionals}
\begin{definition}\ \\
Let $E = \underline{\text{Linear space over } \R \text{ or } \C}$. A linear functional on $E$ is a linear operator $f: E \to \R \text{ or } \C$ s.t. 
$$f(ax + by) = af(x) + bf(y),\ x,y \in E,\ a,b \in \R\text{ or }\C.$$ 
We say $f(\cdot)$ is a \underline{Bounded Linear Functional} if 
$$\norm{f} = \sup_{\norm{x} = 1} \abs{f(x)} < \infty.$$ 
By dilation and additive, properties of $f(\cdot)$ boundedness implies $\abs{f(x - y)} \leq \norm{f} \norm{x - y},\ x,y \in E$. Hence, $f(\cdot)$ is continuous and in fact Lipschitz continuous. Conversely, if a linear functional is continuous then it is bounded. 
\end{definition}
\begin{proof}\ \\
Suppose $f(\cdot)$ is not bounded $\Rightarrow$ $\exists$ sequence $(x_n) \subset E$ s.t. $\abs{f(x_n)} \geq n \norm{x_n},\ n=1,2,\dots$. By linearity, $\abs{f(\frac{x_n}{n \norm{x_n}})} \geq 1,\ n=1,2,\dots$, but $\lim_{n \to \infty} \frac{x_n}{\norm{x_n}} = 0$ and $f(0) = 0$ $\Rightarrow$ $f(\cdot)$ is not continuous at $0$.
\end{proof}

\begin{definition}\ \\
    Let $E$ be a normed space, then space of all \underline{bounded linear functionals} on $E$ is known as the \underline{dual space $E^*$ of $E$}. It is also a normed space with norm $\norm{f} = \sup_{\norm{x} = 1} \abs{f(x)}$, and in fact a Banach space.
\end{definition}

\begin{definition}\ \\
    Let $E$ be a linear space and $H \subset E$ a subspace. We say $H$ is a hyperplane if $\text{codim}(H) = 1$, i.e. $\dim(\frac{E}{H}) = 1$. 
\end{definition}

\np \textbf{\underline{Goal:}} Make an equivalence between \underline{Bounded Linear Functionals} on $E$ and \underline{closed} hyperplanes in $E$. Does there exists a non closed hyperplane? (Not in finite dimensions) Another question is that does there exists a subset $F \subset \R$ which is not Lebesgue
measurable?\\ 
The answers are yes in both cases. However, construction uses axiom of choice.

\begin{proposition}\label{4.4}\ \\
$E$ a linear space,
\begin{enumerate}[label = (\alph*)]
    \item For every linear functional $f$ on $E$, $ker(f)$ is a hyperplane in $E$.
    \item  If $f,g\neq 0$ are linear functionals on $E$ s.t. $ker(f) = ker(g)$, then $f = ag$ for some $a \neq 0$. 
    \item For every hyperplane $H \subset E$, $\exists$ a linear functional $f \neq 0$ on $E$ s.t. $ker(f) = H$.
    \item If $E$ is a Banach space, then $f(\cdot)$ is Bounded \underline{if.f} $ker(f) = H$ is closed.
\end{enumerate}
\end{proposition}
\begin{proof}\ 
\begin{enumerate}[label = (\alph*)]
    \item Let $x,y \notin \ker(f)$ so $f(x), f(y) \neq 0$. Hence $\exists$ scalar $\lambda \neq 0$ s.t. $f(x) = \lambda f(y)$ $\Rightarrow$ $x - \lambda y \in \ker(f)$.\\
    Hence, $[x],[y] \in \frac{E}{\ker(f)}$, then $[x] = \lambda[y] \Rightarrow \dim[\frac{E}{\ker(f)}] = 1$\\
    $f$ is bounded $\Rightarrow$ f continuous $\Rightarrow$ $\ker(f) = f^{-1}(0)$ closed. 
    \item Consider the induced functional $\Tilde{f}, \Tilde{g}: \frac{E}{H} \to \R\text{ or } \C$, $H = \ker(f) = ker(g)$.\\ We have $\dim[\frac{E}{H}] = 1$ \imply $\Tilde{f} = a \Tilde{g}$ for some $a \neq 0$ \imply $f = ag$.
    \item Assume $\dim[\frac{E}{H}] = 1$ \imply $\frac{E}{H} = \{a [x_0]: a \in \C \text{ or } \R\}$ for some $x_0 \in E$.\\
    For $x \in E$, we have $[x] = a(x) [x_0]$ for some $a(x) \in \R \text{ or } \C$. Define $f(x) = a(x)$, then $f$ is linear and $\ker(f) = H$. Assume $E$ a Banach space and $H$ closed, we have $\dim(\frac{E}{H}) = 1$. Recall that $\frac{E}{H}$ is also a Banach space with norm
    $$\norm{[x]} = \inf_{y \in H} \norm{x + y}, x\in E $$
    Let $\Tilde{f}$ be a linear functional on $\frac{E}{H}$. Then we have $\dim(\frac{E}{H})$ finite \imply $\Tilde{f}$ continuous \imply $\abs{\Tilde{f}([x])} \leq A \norm{[x]},\ \forall\ x \in E$. Define $f(x) = \Tilde{f}([x]),\ x \in E$, then $\ker(f) = H$ and $\abs{f(x)} \leq A \norm{[x]} \leq A \norm{x}$.
\end{enumerate}
\begin{remark}
Recall "Every linear mapping on finite-dimensional vector space is continuous".
\end{remark}
\end{proof}



\vspace{12pt}
\subsection{Riesz Representation Theorem}

\begin{theorem}[Riesz Representation Theorem]\label{RRT}\ \\
Let $\Hs$ be a Hilbert space.
\begin{enumerate}[label = (\alph*)]
    \item For every $y \in \Hs$, the fucntion $f(x) = \inprod{x}{y},\ x \in \Hs$ is a bounded linear functional on $\Hs$.
    \item If $f: \Hs \to \C \text{ or } \R$ is a bounded linear functional on $\Hs$, then $\exists\ y \in \Hs$ s.t. $f(x) = \inprod{x}{y},\ x\in \Hs$. Hence the dual $\Hs^*$ of $\Hs$ is isometric to $\Hs$.
\end{enumerate}
\end{theorem}
\begin{proof}\
\begin{enumerate}[label = (\alph*)]
    \item It follows from Cauchy-Schwarz s.t.
    \begin{equation*}
        \abs{f(x)} = \abs{\inprod{x}{y}} \leq \norm{x} \norm{y} 
    \end{equation*}
    and we have $\norm{f} = \norm{y}$ by setting $x = \frac{y}{\norm{y}}$. 
    \begin{remark}\ \\
    In this case, the supremum of norm is attained by $\norm{f} = \sup_{\norm{x} = 1} \abs{f(x)} = f(\frac{y}{\norm{y}})$. It is also true for any \underline{finite} dimensional space, since unit ball is compact. In general it is not true for $\infty$-dimensional Banach spaces where the unit ball is not compact.
    \end{remark}
    \item Let $f: \Hs \to \R \text{ or } \C$ be a bounded linear functional on $\Hs$. Let $H = \ker(f)$, and $H$ is closed.\\ Let $H^\perp$ be orthogonal complement of $H$ s.t. $\Hs = H \bigoplus H^\perp$. We have $\dim[\frac{\Hs}{H}] = 1$ \imply $\dim(H^\perp) = 1$.\\
    Therefore, we can write
    \begin{equation*}
        \Hs = H + \vspan{x_0},\ \text{ for some } y_0 \in H^\perp.
    \end{equation*}
    Consider the map
    \begin{equation*}
        g(x) := \inprod{x}{y_0},
    \end{equation*}
    by part (a) we have
    \begin{equation*}
        g \in \Hs^*,\ \ker{g} = \ker{f}.
    \end{equation*}
    Therefore, by \hyperref[4.4]{Proposition 4.4}, we have $f = ag$ for some constant $a$, and it follows that
    \begin{equation*}
        f(x) = ag(x) = \inprod{x}{a y_0}.
    \end{equation*}
    Complete the proof.
\end{enumerate}
\end{proof}


\begin{theorem}[Radon-Nikodym Theorem]\label{R-N thm}\ \\
Let $\mu,\nu$ be two $\sigma$-finite measures s.t. $\nu \ll \mu$, i.e. $\nu$ is \underline{absolutely continuous} w.r.t. $\mu$. (i.e. $\mu(A) = 0$ \imply $\nu(A) = 0$) Then $\exists\ g \geq 0$ s.t. $g$ is $\mu$-integrable and $\nu(A) = \int_A g d\mu$ for measurable $A$, where $g$ is referred to as R-N Derivative $g = \frac{d\nu}{d \mu}$.
\end{theorem}
\begin{proof}\ \\
Consider the linear functional $F: \Ls^2(\mu) \to \R \text{ or } \C$ s.t. 
\begin{equation*}
    F(f) = \int_\Omega f d\mu.
\end{equation*}
Then we have $\norm{F(f)} \leq \norm{f}_2 \mu(\Omega)^{\frac{1}{2}}$ by H\"older inequality. $F$ is also a bounded linear functional on $\Ls^2(\mu + \nu)$, hence, the \hyperref[RRT]{Riesz Representation Theorem} implies $\exists\ h \in \Ls^2(\mu+\nu)$ s.t.
$$F(f) = \int f h d(\mu+\nu)$$ 
for $f \in \Ls^2(\mu + \nu)$, i.e. 
\begin{equation}
    \int_\Omega f d\mu = \int_\Omega fh d\mu + \int_\Omega fh d\nu.\label{eq:4.1}
\end{equation} 
Rearranging terms, we have
\begin{equation}
    \int_\Omega fh d\nu = \int_\Omega f(1 - h) d\mu.\label{eq:4.2}
\end{equation} 
Let $A \subset \Omega$, $A$ measurable, set $f = \frac{1}{h} \1_A$, we have $\nu(A) = \int_A  g d\mu,\ g = \frac{1-h}{h}$ \imply $g = \frac{d\nu}{ d \mu}$. It remains to show this is well-defined, i.e. $h \neq 0$.
\begin{claim}
$0 < h \leq 1$ a.e. \\
Note $\mu(A) = 0 \Leftrightarrow \mu(A) + \nu(A) = 0$. Let $A = \{h \leq 0\},\ f = \1_A$. \hyperref[eq:3.1]{Equation 3.1} implies that we have 
\begin{equation*}
    \int_A h [d\mu + d\nu] \leq 0 \Rightarrow \mu(A) = 0 \Rightarrow h > 0\ \mu-a.e.
\end{equation*}\\
To show $h \leq 1$, set $B = \{h > 1\}$, and let $f = \1_B$. \hyperref[eq:4.1]{Equation 4.1} implies 
\begin{equation*}
    \mu(B) = \int_B h [d\mu + d\nu] > \mu(B) \text{ unless } \mu(B) = 0
\end{equation*}
which shows that $h \leq 1$.
\end{claim}
Now that we know $0 < h \leq 1$, we have \hyperref[eq:4.2]{Equation 4.2} holds for all $f \geq 0,\ f \in \Ls^2(\mu + \nu)$. By using Monotone Convergence Theorem, we can conclude that \hyperref[eq:4.2]{Equation 4.2} holds for all $f \geq 0$ (Thus not restrict on finite measure $\mu,\nu$, but extend to $\sigma$-finite), and both sides could be $+\infty$. Now can plug in $f = \frac{1}{h}\1_A$ to get the result.\\ 
\end{proof}
\begin{remark}\ \\
R-N derivative $\frac{d\nu}{d\mu}$ is \underline{unique} since $\int_A g d\mu = 0$ for all $\mu$-measurable $A$ \imply $g = 0$ $\mu$-a.e
\end{remark}

\begin{remark}\ \\
For $\Hs = \text{Hilbert space}$, \hyperref[RRT]{Riesz Representation Theorem} identifies the dual space $\Hs^*$ and can imply R-N Theorem.
\end{remark}

\begin{corollary}\ \\   
Consider spaces $\Ls^p(\Omega, \mu)$ for
$1 \leq p < \infty$. R-N Theorem implies dual of $\Ls^p(\Omega, \mu)$ is $\Ls^{p'}(\Omega, \mu)$, where $\frac{1}{p} + \frac{1}{p'} = 1$ (also for $p=1,\ p' = +\infty$).
\end{corollary}
\begin{proof}\
\begin{itemize}
    \item Easy part:\\
    $g \in \Ls^{p'}$ induces a bounded linear functional on $\Ls^p$ (\hyperref[RRT]{Riesz Representation Theorem}) by setting 
    \begin{equation*}
        F(f) = \int_\Omega fg d\mu.
    \end{equation*} 
    Use H\"older inequality, we get
    \begin{equation*}
        \abs{F(f)} \leq \norm{f}_p \norm{g}_{p'},\ \norm{F} = \norm{g}_{p'}    
    \end{equation*} 
    by choosing $f = \abs{g}^{p' - 1} \frac{\bar{g}}{\abs{g}}$, since 
    \begin{equation*}
        F(f) = \int_\Omega \abs{g}^{p'} d\mu = \norm{g}_{p'}^{p'}
    \end{equation*}
    From $\frac{1}{p} + \frac{1}{p'} = 1$ \imply $p' - 1 = \frac{p'}{p}$, we have 
    \begin{equation*}
        \norm{f}_{p}^p = \int \abs{f}^p d\mu = \int \abs{g}^{p'} d\mu = \norm{g}_{p'}^{p'}\ \Rightarrow\ \norm{f}_p = \norm{g}_{p'}^{p'/p} = \norm{g}_{p'}^{p' - 1}
    \end{equation*}
    which implies
    \begin{equation*}
      \norm{g}_{p'}^{p'} = \norm{g}_{p'} \norm{f}_p^p\ \Rightarrow\ \norm{F} = \norm{g}_{p'}.  
    \end{equation*}
    \begin{remark}
    Note $\sup_{\norm{f} = 1} \abs{F(f)}$ is attained for $1 < p < \infty$ with $f = g^{p'-1} \frac{\bar{g}}{\abs{g}}$, which shows $F$ is bounded and 
    $$\Ls^{p'}(\Omega,\Sigma, \mu) \subset \text{ dual of } \Ls^p(\Omega, \Sigma, \mu).$$ The supremum is also attained when $p = \infty$, $g \in \Ls^1$ and let $f = \frac{\bar{g}}{\abs{g}}$. The supremum is not attained when $p = 1$, $g \in \Ls^\infty$. Let 
    \begin{equation*}
        F(f) = \int fg d\mu.
    \end{equation*}
    If $g(\cdot)$ is continuous on $\R$ with unique maximum, then $\sup_{\norm{f} = 1} \abs{F(f)}$ is not attained.  
    \end{remark}
    \begin{proof}[Summary]\ \\
    If $1 \leq p \leq \infty$, $\Ls^{p'}$ is contained in the dual of $\Ls^p$.\\ 
    If $1 < p \leq \infty$, then $\sup_{\norm{f}_p = 1} \abs{F(f)}$ is attained.\\ 
    For $p = 1$, the supremum is not attained.
    \end{proof}
    \item For $E = \Ls^p(\Omega,\Sigma,\mu)$ with $1 \leq p < \infty$, and $F \in E^*$, we consider finite measure space $\mu(\Omega) < \infty$. Define measure $\nu$ on $\Sigma$ by
    \begin{equation*}
        \nu(A) = F(\1_A) = \int_A h \dr \mu,\ A \in \Sigma,
    \end{equation*}
    where $h \in \Ls_0(\mu)$ by the \hyperref[RRT]{Riesz Representation Theorem}, and we have
    \begin{center}
        $\mu(A) = 0$ \imply $\nu(A) = 0$ \imply $\nu \ll \mu$.
    \end{center}
    By \hyperref[R-N thm]{Radon-Nikodym Theorem}, we know
    \begin{equation*}
        \nu(A) = \int_A g d\mu
    \end{equation*}
    for some $g \in \Ls^1$ s.t. $g = \frac{d\nu}{d\mu}$. It suffices to show $g \in \Ls^{p'}$. Note that $p' > 1$, we have
    \begin{equation*}
        F(f) = \int_{\Omega} fg d\mu
    \end{equation*}
    for simple functions $f$. We can decompose $f = f^+ - f^-,\ g = g^+ - g^-$, find $(f_n^{\pm} \in \Ls^p)$ s.t. $f_n^{\pm} \uparrow f^{\pm}$, and apply MCT to four different cases, and thus we have
    \begin{equation*}
            F(f) = \int_{\Omega} fg d\mu
    \end{equation*}
    holds for all $f \in \Ls^p$. Can set $f = g^{p' - 1}$, and use 
    \begin{equation}
        \abs{F(f)} \leq \norm{F} \norm{f}_p\label{eq:4.3}
    \end{equation}
    where $\frac{1}{p} + \frac{1}{p'} = 1$ \imply $p' - 1 = \frac{p'}{p}$. And \hyperref[eq:4.3]{Equation 4.3} implies
    \begin{equation*}
        \int g^{p'} d\mu \leq \norm{F} (\int g^{p'} d\mu)^{\frac{1}{p}},
    \end{equation*}
    and thus
    \begin{align*}
        \norm{g}_{p'}^{p'} &\leq \norm{F} \norm{g}_{p'}^{\frac{p'}{p}}\\
        &= \norm{F} \norm{g}_{p'}^{p' - 1}
    \end{align*}
    which concludes $\norm{g}_{p'} \leq \norm{F} < +\infty$. Complete the proof.
\end{itemize}
\end{proof}

\begin{remark}\ \\
$\Ls^1$ is a subset of the dual of $\Ls^\infty$, but not equal to it. Note if $F: \Ls^\infty (\mu) \to \C$ is a bounded linear functional, then if $\Omega = \mathrm{K} = \text{ compact Hausdorff space } F$ induces a bounded linear functional on $C(K) = \text{ space of continuous functions on } K$, we have $C(k) \subset \Ls^\infty (K, \Sigma, \mu)$, where $\Sigma = \Bs(K)$. \textcolor{red}{The counterexample is a little complex, check later}
\end{remark}

\begin{theorem}[Riesz Representation Theorem, v.2]\label{RRT2}\ \\
Let $E = C(K)$ be the space of continuous functions on compact Hausdorff space $K$. Then,
\begin{enumerate}[label = (\alph*)]
    \item For every Borel regular signed measure on $K$, the functional $F(f) = \int_K f d\mu$ is a bounded linear functional on $K$.
    \item Every bounded linear functional on $C(K)$ can be expressed in (a) for some measurable $\mu$, and $\norm{F} = \abs{\mu} (K)$, i.e. $\text{Tv}(K)$, where $\text{Tv}(\cdot)$ denotes the total variation. 
\end{enumerate}
\end{theorem}


\vspace{12pt}
\subsection{Hahn-Banach Theorem}
\begin{theorem}[Hahn-Banach Theorem]\label{HB thm}\ \\
Let $E_0$ be a subspace of a normed space, then every bounded linear functional $f_0: E_0 \to \R \text{ or } \C$ has an extension $f: E \to \R \text{ or } \C$ s.t. $\norm{f} = \norm{f_0}$.
\end{theorem}
\begin{proof}\ \\
Assume $E$ is separable, otherwise we need "transfinite induction". 
\begin{remark}
Separability allows us to extend $f_0$ 1 dimension at a time.
\end{remark}
Let $\{x_n: n = 1,2,\dots\}$ have the property that its span is dense in $E$. We can do induction:
\begin{equation*}
    E_0 \to E_0 + \{x_1\} \to E_0 + \{x_1,x_2\} \to \dots \to E_0 + \vspan{x_1,\dots,x_n},\ \norm{f} = \norm{f_0}.
\end{equation*}
This final space is dense in $E$ and so can extend $f$. Thus it suffices to prove the extension by 1 dimension s.t. $E \to E_0 + \{x_1\}$.

\np Note the extension is determined by a single number $\gamma = f(x_1)$. First we consider linear functional $f_0: E_0 \to \R$ and extend to $f: E_0 + \{x_1\} \to \R$ with $\norm{f} = \norm{f_0}$, which is equivalent to (WLOG, assuming $\norm{f} = \norm{f_0} = 1$)
\begin{equation*}
    \abs{f_0(x_0) + \lambda \gamma} \leq \norm{x_0 + \lambda x_1},\ \forall\ x_0 \in E_0,\ \lambda \in \R.
\end{equation*}
Divide inequality by $\lambda \neq 0$, it suffices to find $\gamma$ s.t.
\begin{equation*}
    \abs{f_0(x_0) + \gamma} \leq \norm{x_0 + x_1},\ \forall\ x_0 \in E_0.
\end{equation*}
Note we have $f_0$ is a real functional, it is equivalent to
\begin{equation*}
    - \norm{x_0 + x_1} \leq f_0(x_0) + \gamma \leq \norm{x_0 + x_1},\ \forall\ x_0 \in A.
\end{equation*}
Such a $\gamma$ exists provided
\begin{equation*}
    \norm{x_0 + x_1} - f_0(x_0) \geq - \norm{x_0'+ x_1} - f(x_0'),\ \forall\ x_0,x_0' \in E_0,
\end{equation*}
which is equivalent to (assuming $\norm{f_0} = 1$)
\begin{align*}
    f_0(x_0 - x_0') &\leq \norm{x_0 + x_1} + \norm{x_0' + x_1},\ \forall\ x_0,x_0' \in E_0\\
    f_0(x_0 - x_0') &\leq \norm{x_0 + x_1} + \norm{-x_0' - x_1},\ \forall\ x_0,x_0' \in E_0
\end{align*}
which can be given by
\begin{equation*}
    f_0(x_0 - x_0') \leq \norm{x_0 - x_0'} \underbrace{\leq}_{\text{triangle inequality}} \norm{x_0 + x_1} + \norm{-x_1 - x_0'}
\end{equation*}

\np \underline{Extension to linear functionals over complexs}:\\
Let $f: E \to \C$ be a linear functional over $\C$, set $g(x) = \Re{f(x)}$, then $g: E \to \R$ is a real linear functional, and $f(x) = g(x) - i g(ix),\ x \in E$. Note
\begin{align*}
    f(ix) &= i f(x)\\
    g(ix) &= - \Im{f(x)}
\end{align*}
Conversely, if $g: E \to \R$ is a real linear functional on Banach space (not necessarily need Banach) $E$ over $\C$, then
\begin{equation*}
    f : E \to \C
\end{equation*}
defined by
\begin{equation*}
    f(x) = g(x) - i g(ix),\ x \in E
\end{equation*}
is a complex linear functional on $E$.

\np To extend $f_0: E_0 \to \C$, we extend $g_0 = \Re{f_0}$ 2 real dimensions s.t.
\begin{equation*}
    E_0 \to E_0 + x_1 \to E_0 + x_1 + i x_1
\end{equation*}
and define $f(\cdot) = g(\cdot) - i g(i \cdot)$.

\np To show $\norm{f} = \norm{f_0}$, we use that fact that for $x \in E_0 + \{\lambda x_1:\ \lambda \in \C\}$, we have
\begin{equation*}
    e^{i \theta} f(x) = f(x e^{i\theta}),\ \theta \in \R.
\end{equation*}
We can choose $\theta$ s.t. $f(x e^{i\theta}) = \Re{g(x e^{i \theta})}$. Now we already have
\begin{equation*}
    \abs{g(x e^{i \theta})} \leq \norm{f_0} \norm{x e^{i \theta}}
\end{equation*}
which implies 
\begin{equation*}
    \abs{f(x)} \leq \norm{f_0} \norm{x},\ \forall\ x \in E_0 + \{\lambda x_1:\ \lambda \in \C\}.
\end{equation*}
\end{proof}







\vspace{12pt}
\subsubsection{Implications of the Hahn-Banach Theorem}
\begin{proposition}[Supporting Hyperplane Theorem]\label{SHT}\ \\
Let $E$ be a normed space, for every $x \in E$, $\exists\ f \in E^*$ s.t. 
\begin{equation*}
    \norm{f} = 1,\ f(x) = \norm{x},
\end{equation*}
i.e. $\sup_{\norm{y}=1} \abs{f(y)}$ is attained at $y = x$. 
\end{proposition}
\begin{proof}\ \\
Consider 1-dimensional space $E_0 = \vspan{x} = \{tx: t \in \R \text{ or } \C\}$. Define $f_0: E_0^* \to \R \text{ or } \C$ s.t.
\begin{equation*}
    f_0(tx) = t \norm{x}.
\end{equation*}
We have $\norm{f_0} = 1$, and by \hyperref[HB thm]{Hahn-Banach Theorem}, $\exists\ f \in E^*$ s.t. $\ f(x) = \norm{x},\ \norm{f} = 1$.
\end{proof}
\begin{remark}[Geometric Interpretation]\ \\
Suppose $B = \text{ unit ball } \{x \in E: \norm{x} \leq 1\}$. Choose $x_0 \in \partial B,\ \norm{x_0} = 1$. $\exists\ f \in E^*$ s.t. $\norm{f} = 1,\ f(x) = \norm{x}$. Let 
\begin{equation*}
    H = \ker{f} + x_0,
\end{equation*}
where $\ker{f} = \{x: f(x) = 1\}$. $H$ intersects $E$ into two disjoint subsets, and $B$ lies entirely in one of them, i.e. for all $x \in B$ we have
\begin{center}
    $\norm{x} < 1$ \imply $\abs{f(x)} \leq \norm{f} < 1$ \imply f(x) < 1,
\end{center}
and we also have $E = \{x: f(x) < 1\} \kup H \kup \{x: f(x) < 1\}$.
\end{remark}
\begin{remark}\ \\
Tangent hyperplane $H$ is not necessarily unique. We can extend this to prove the existence supporting hyperplane for more general convex sets.
\end{remark}

\vspace{12pt}
\subsection{Second Duality}

\begin{definition}\ \\
    Let $E$ be a Banach space, the \underline{duality of $E$} is the space of all bounded linear functionals on $E$. We have for $f \in E^*,\ \norm{f} = \sup_{\norm{x} = 1} \abs{f(x)}$, and $E^*$ is a Banach space.  
\end{definition}

\np Let $E^*$ be the dual of $E$, and $E^{**}$ be the dual of $E^*$. $\exists$ natural embedding $E \to E^{**}$, $x \to x^{**} \in E^{**}$ s.t.
\begin{equation*}
    x^{**}(f) = f(x),\ f \in E^*.
\end{equation*}
We have
\begin{equation*}
    \norm{x^{**}} = \sup_{\norm{f} = 1,\ f \in E^*} \abs{x^{**} (f)} = \sup_{f \in E^*,\ \norm{f} = 1} \abs{f(x)} \leq \norm{x},
\end{equation*}
which implies
\begin{equation*}
    \norm{x^{**}} \leq \norm{x},\ x \in E.
\end{equation*}
By Hahn-Banach Theorem, we have $\norm{x^{**}} = \norm{x}$. Recall 
\begin{equation*}
    f_0: f_0(tx) = t \norm{x},\ t \in R \text{ or } \C,\ \norm{f_0} = 1.
\end{equation*}
Extend $f_0$ to a functional $f_x: E \to \R \text{ or } \C$ s.t. $\norm{f_x} = 1,\ f_x(x) = \norm{x}$ \imply $x^{**}(f_x) = \norm{x}$ \imply $\norm{x^{**}} \geq \norm{x}$.

\np Conclude the embedding $E \to E^{**}$ is an isometry if the mapping is onto, say $E$ is a reference space.

\begin{theorem}[Second Dual Space]\ \\
Let $E$ be a normed space, then $E$ can be considered as a linear subspace of $X^{**}$. For this, a vector $x \in E$ is considered as a bounded linear functional on $E^*$ via the action:
\begin{equation*}
    x: f \to f(x),\ f \in E^*.
\end{equation*}
\end{theorem}


\begin{examples}\ 
\begin{enumerate}[label = (\alph*)]
    \item Hilbert spaces.
    \item $E = \Ls^p$ spaces for $1 < p < \infty$, we have
    \begin{equation*}
        E^* = \Ls^{p'},\ \frac{1}{p} + \frac{1}{p'} = 1,\ 1 < p' < \infty
    \end{equation*}
    and $E^{**} = \Ls^p$.
\end{enumerate}
\end{examples}


\begin{definition}[Reflexivity]\ \\
    A normed vector space $E$ is called \textit{reflexive} if $E^{**} = E$ under the natural embedding.
\end{definition}

\begin{proposition}\ \\
If $E$ is a reflexive normed space, then every functional $f^* \in E^*$ attains its maximum on $E$.
\end{proposition}
\begin{proof}\ \\
E is reflexive and $f \in E^*$, then $\exists\ x_f \in E^{**} = E$ with $\norm{x_f} = 1$ and $\norm{f} = f(x_f)$, i.e. $\sup_{\norm{x} = 1} \abs{f(x)}$ is achieved at $x = x_f$. This follows from \hyperref[SHT]{Supporting Hyperplane Theorem}.
\end{proof}

\np To show a normed space $E$ is not reflexive, it suffices to find $f \in E^*$ s.t. $\sup_{\norm{x} = 1} \abs{f(x)}$ is note attained. 

\begin{example}\ \\
Let $E = C([0,1]) = \text{ space of continuous functions } g: [0,1] \to \C$, and $\norm{g} = \sup_{0 \leq t\leq 1} \abs{g(t)}$. For $f \in E^*$, we define
\begin{equation*}
    f(g) = \int_0^1 h(x) g(x) dx,
\end{equation*}
where $h(x) = -1$ for $0  < x < \frac{1}{2}$, and $h(x) = 1$ for $\frac{1}{2} < x < 1$. Then
\begin{equation*}
    \norm{f} = 1 = \sup_{\norm{g} = 1} \abs{f(g)},
\end{equation*}
but sup is not attained.
\end{example}

\subsection{Separating Hyperplane Theorem}
\begin{remark}
Separating Hyperplane Theorem is an extension of \hyperref[SHT]{Supporting Hyperplane Theorem}, generalizing to arbitrary convex sets.
\end{remark}

\subsubsection{Sublinear Functionals}
\begin{definition}\ \\
    Let $E$ be a linear space, a function
    \begin{equation*}
        \norm{\cdot} : E \to [0, \infty)
    \end{equation*}
    is \textit{sublinear} if it satisfies
    \begin{enumerate}[label = (\alph*)]
        \item $\norm{\lambda x} = \lambda \norm{x},\ \lambda \in \R,\ \lambda > 0,\ x \in E$.
        \item $\norm{x + y} \leq \norm{x} + \norm{y},\ x,y \in E$.
    \end{enumerate}
\end{definition}
\begin{remark}\ \\
Note the difference of the definition of sublinear functions and the norm. For a sublinear function $\norm{\cdot}$, it requires two additional properties to become a norm:
\begin{itemize}
    \item $\norm{-x} = \norm{x},\ \forall\ x\in E$.
    \item $\norm{x} = 0 \Rightarrow x = 0$.
\end{itemize}
\end{remark}

\vspace{6pt}
\begin{theorem}[Extension of Hahn-Banach Theorem]\ \\
Let $E_0$ be a subspace of a linear vector space over $\R$, and $\norm{\cdot}$ be a sublinear functional on $E$ and $f_0: E_0 \to \R$ a linear functional on $E_0$ satisfying
\begin{equation*}
    f_0(x) \leq \norm{x},\ x \in E_0.
\end{equation*}
Then $f_0$ admits an extension $f$ on $E$ s.t.
\begin{equation*}
    f(x) \leq \norm{x},\ x \in E.
\end{equation*}
\end{theorem}
\begin{proof}\ \\
Same argument as in the proof of \hyperref[HB thm]{Hahn Banache Theorem}.
\end{proof}
\begin{remark}\ \\
There is no guarantee that if $E$ is a Banach space, then $f(\cdot)$ is a bounded linear functional.
\end{remark}

\subsubsection{Geometric Properties of Sublinear Functionals}
\begin{definition}\ \\
    A subset $K$ of a linear vector space $E$ is \textit{absorbing} if
    \begin{equation*}
        E = \bigcup_{t \geq 0} t K,
    \end{equation*}
    where $tK = \{tk: k \in K\}$.
\end{definition}

\begin{proposition}[Minkowski Functional]\ \\
Let $K \subset E$ be an absorbing convex subset of a linear vector space $E$ s.t. $0 \in K$. Then the \textit{Minkowski functional}
\begin{equation*}
    \norm{x}_K = \inf\{t > 0: x \in tK\}
\end{equation*}
is a sublinear functional on $E$. Conversely, let $\norm{\cdot}$ be a sublinear functional on $E$. Then the \textit{sub-level} set
\begin{equation*}
    K = \{x \in E:\ \norm{x} \leq 1\}
\end{equation*}
is an absorbing convex set and $0 \in K$.
\end{proposition}
\begin{proof}\ \\
\begin{itemize}
    \item $\Rightarrow:$\\
    Let $K$ be an absorbing convex set with $0 \in K$, the main observation is that since $0 \in K$, and $K$ is convex, we have
    \begin{equation*}
        x \in K \Rightarrow tx \in K,\ 0 \leq t < 1.
    \end{equation*}
    WTS $\norm{\lambda x} = \lambda \norm{x},\ x \in E,\ \lambda > 0$. Note
    \begin{align*}
        \norm{\lambda x} &= \inf\{t > 0:\ \lambda x \in tK\}\\
        &= \inf\{t > 0:\ x \in \frac{t}{\lambda} K\}\\
        &= \lambda \inf\{s > 0:\ x \in sK\}\\
        &= \lambda \norm{x}.
    \end{align*}
    It remains to prove the triangle inequality. For $x\in tK,\ y \in sK$, we have 
    \begin{equation*}
        x = tk_1,\ y = sk_2,\ k_1,k_2 \in K,
    \end{equation*}
which implies
\begin{align*}
    x + y &= (t + s) [\frac{t}{t+s} k_1 + \frac{s}{t+s} k_2]\\
    &= (t+s) k
\end{align*}
for some $k \in K$, since $K$ is convex. Therefore, we have
\begin{equation*}
    x + y \in (t+s)K \Rightarrow \norm{x + y} \leq \norm{x} + \norm{y},
\end{equation*}
since $\{t: x \in tK\}$ is an interval $(\alpha(x), \infty)$.
\item $\Leftarrow:$\\
$0 \in K$ since $\norm{0} = 0$, and the convexity comes from triangle inequality.\\
\begin{remark}
If $K \neq -K$, then $\exists\ x \in E$ with $\norm{x} \neq \norm{-x}$ if $K = E$, then $\norm{\cdot} = 0$.
\end{remark}
\end{itemize}
\end{proof}


\subsubsection{Seperating Hyperplane}

\begin{theorem}[Separation of a Point from a Convex Set]\ \\
Let $K$ be an open convex subset of a normed space $E$ and $x_0 \neq K$. Then exists a continuous linear functional
\begin{equation*}
    f: E \to \R,\ f \not\equiv 0,
\end{equation*}
and $f(x) < f(x_0),\ x \in K$.
\end{theorem}
\begin{proof}\ \\
By translation can assume WLOG $0 \in K$. Since $K$ is open, it is absorbing. Let $\norm{\cdot}_K$ be Minkowski functional, then
\begin{equation*}
    \norm{x}_K \leq \frac{1}{r} \norm{x},\ x \in E
\end{equation*}
if $B(0,r) \subset K$. Proceed as in separation theorem for a unit ball. Define $f_0$ on $\vspan{x_0}$ by
\begin{equation*}
    f(tx_0) = t\norm{x_0},\ t \in \R,
\end{equation*}
then if $E_0 = \{\lambda x_0:\ \lambda \in \R\}$, we have
\begin{equation*}
    f_0(x) \leq \norm{x}_K,\ x \in E_0,
\end{equation*}
since \begin{equation*}
    \norm{tx}_K = t \norm{x}_K,\ t \geq 0.
\end{equation*}
\hyperref[HB thm]{Hahn Banach Theorem} implies that we can extend $f_0$ to $f: E \to \R$ s.t.
\begin{equation*}
    f(x) \leq \norm{x}_K \leq \frac{1}{r} \norm{x},\ x \in E,
\end{equation*}
and thus $f$ is a bounded linear functional.\\
For separation, we have
\begin{equation*}
    f(x) \leq \norm{x}_K \leq 1 \leq \norm{x_0}_K = f_0(x_0) = f(x_0),\ x \in K,
\end{equation*}
which implies $f(x) \leq f(x_0),\ x \in K$, and since $K$ is open,
\begin{equation*}
    x + tv \in K
\end{equation*}
for some $t > 0$, $\forall\ v$ s.t. $\norm{v} = 1$. Hence,
\begin{equation*}
    f(x + tv) \leq f(x_0),\ t = t_x > 0,\ \norm{x} = 1,
\end{equation*}
and thus 
\begin{equation*}
    f(x) + t \sup_{\norm{v} = 1} f(v) \leq f(x_0).
\end{equation*}
Note that $\sup_{\norm{v} = 1} f(v) > 0$, we have
\begin{equation*}
    f(x) < f(x_0).
\end{equation*}
\end{proof}



\begin{theorem}[Separation of Convex Sets]\label{SepConvSets}\ \\
Let $A,B$ be disjoint convex subsets of a normed space $E$,
\begin{enumerate}[label = (\alph*)]
    \item If $A$ is open, then $\exists$ a bounded linear functional
    \begin{equation*}
        f: E \to \R
    \end{equation*}
    s.t. $f(a) < f(b),\ \forall\ a \in A,\ b \in B$.
    \item If $A$ is closed and $B$ is compact, then
    \begin{equation*}
        \sup_{a \in A} f(a) < \inf_{b \in B} f(b)
    \end{equation*}
\end{enumerate}
\end{theorem}
\begin{proof}\
\begin{enumerate}[label = (\alph*)]
    \item Let 
    \begin{equation*}
        K = A -B = \{a - b:\ a \in A,\ b \in B\},
    \end{equation*}
    then $K$ is open, convex and $0 \in K$. By the previous theorem, $\exists\ f \in E$ s.t.
    \begin{equation*}
        f(a - b) < f(0) = 0,\ \forall\ a \in A,\ b \in B.
    \end{equation*}
    Therefore,
    \begin{equation*}
        f(a) < f(b),\ a \in A,\ b \in B.
    \end{equation*}
    \item Suppose $A$ closed and $B$ compact, we have
    \begin{equation*}
        d(A,B) := \inf\{\norm{x - y}:\ x \in A,\ y \in B\} = \delta > 0.
    \end{equation*}
    Define $A_\delta = \{x \in E:\ d(x,A) < \delta\}$, we have $A_\delta$ is open, and take
    \begin{equation*}
        \ep = \frac{\delta}{2},
    \end{equation*}
    we have $A_\ep \cap B = 0$. By part (a), we have $\exists\ f \in E^*$ s.t.
    \begin{equation*}
        f(x) < f(y),\ \forall\ x \in A,\ y \in B.
    \end{equation*}
    For $a \in A$, we have $a + \frac{\delta}{2} v \in A_\delta$ if $\norm{v} = 1$. Thus
    \begin{equation*}
        f(a + \frac{\delta}{2} v) = f(a) + \frac{\delta}{2} f(v) < f(b),\ b \in B.
    \end{equation*}
    Take sup over $\norm{v} = 1$, we have
    \begin{equation*}
        \sup_{\norm{v} = 1} \abs{f(v)} = \ep > 0,
    \end{equation*}
    which implies 
    \begin{equation*}
        f(a) < f(b) - \ep,\ a \in A,\ b \in B,
    \end{equation*}
    which further implies
    \begin{equation*}
        \sup_{a \in A} f(a) < \inf_{ b \in B} f(b).
    \end{equation*}
\end{enumerate}
\end{proof}

\begin{corollary}\ \\
Let $K \subset E$ be a closed convex set, then $K$ is the intersection of all half spaces containing $K$, where half space $H = \{x \in E:\ f(x) \leq \lambda,\ f \in E^*\}$.
\end{corollary}
\begin{proof}\ \\
If $x \notin K$, then $\exists\ f \in E^*$ s.t.
\begin{equation*}
    \sup_{k \in K} f(k) < f(x_0).
\end{equation*}
Can choose $\lambda$ s.t. $\sup_{k \in K} f(k) \leq \lambda \leq f(x_0)$.
\end{proof}

\vspace{12pt}
\subsection{Bounded Linear Operators}

\begin{definition}\ \\
    $X,Y$ Banach spaces, we say $T: X \to Y$ is a \textit{linear operator} if 
    \begin{equation*}
        T(ax + by) = a Tx + b Ty,\ a,b \in \R \text{ or }\C,\ x,y \in X.
    \end{equation*}
    The operator $T$ is said to be \textit{bounded} if
    \begin{equation*}
        \norm{T} := \sup_{\norm{x} = 1} \norm{Tx} < \infty,
    \end{equation*}
    then $T$ is \textit{Lipschitz}:
    \begin{equation*}
        \norm{Tx - Ty} \leq \norm{T} \norm{x - y},\ x,y \in X.
    \end{equation*}
\end{definition}

\begin{remark}\ \\
The operator norm is a norm on bounded linear operators.
\end{remark}

\vspace{6pt}
\begin{proposition}\ \\
Let $\Ls(X,Y) = \text{ space of bounded linear operators } T: X \to Y$, then $\Ls(X,Y)$ is a \textit{Banach space} with the norm $T \to \norm{T}$, i.e.
\begin{enumerate}
    \item $\norm{T} = 0 \Leftrightarrow T = 0$.
    \item $\norm{\lambda T } = \abs{\lambda} \norm{T},\ \lambda \in \R \text{ or } \C,\ T \in \Ls(X,Y)$.
    \item $\norm{T + S} \leq \norm{T} + \norm{S}$, $T,S \in \Ls(X,Y)$.
    \item In addition operator norm satisfies
    \begin{equation*}
        \norm{TS} \leq \norm{T} \norm{S},\ T,S \in \Ls(X,Y).
    \end{equation*}
\end{enumerate}
\end{proposition}

\vspace{6pt}
\begin{definition}[Adjoint operator]\ \\
Suppose $T \in \Ls(X,Y)$, then the adjoint operator $T^*$ of $T$ is defined as (for $f \in Y^*$)
\begin{equation*}
    T^* f : X \to \R \text{ or } \C,\ T^*f(x) = f(Tx),
\end{equation*}
and we further have
\begin{itemize}
    \item $\abs{T^* f(x)} = \abs{f(Tx)} \leq \norm{f} \norm{Tx} \leq \norm{f} \norm{T} \norm{x}$
    \item $T^* f$ is a linear functional.
    \item $\norm{T^* f} = \sup_{\norm{x} = 1} \norm{f}\norm{T}\norm{x} = \norm{f} \norm{T}$.
\end{itemize}
Hence $T^* f \in X^*$, and 
\begin{equation*}
    \norm{T^* f} \leq \norm{T}\norm{f},
\end{equation*}
and thus $T^*: Y^* \to X^*$ is a linear operator and bounded with $\norm{T^*} \leq \norm{T}$.
\end{definition}



\begin{proposition}\ \\
For every $T \in \Ls(X,Y)$, the adjoint $T^*$ is in $\Ls(Y^*, X^*)$ and $\norm{T^*} = \norm{T}.$
\end{proposition}
\begin{proof}\ 
\begin{align*}
    \norm{T^*} &= \sup_{\norm{f}_{Y^*} = 1} \norm{T^* f}_{X^*}\\ 
    &= \sup_{\norm{f}_{Y^*} = 1} \sup_{\norm{x}_X = 1} \abs{T^* f(x)}\\
    &= \sup_{\norm{f}_{Y^*} = 1} \sup_{\norm{x}_X = 1} \abs{f(Tx)}\\
    &=  \sup_{\norm{x}_X = 1} \sup_{\norm{f}_{Y^*} = 1}\abs{f(Tx)}.
\end{align*}
From \hyperref[HB thm]{Hahn Banach Theorem}, 
\begin{equation*}
    \sup_{\norm{f}_{Y^*} = 1} \abs{f(Tx)} = \norm{Tx},
\end{equation*}
conclude that
\begin{equation*}
    \norm{T^*} = \sup_{{x}_X = 1} \norm{Tx} = \norm{T}.
\end{equation*}
\end{proof}


\begin{proposition}[Properties of Adjoint Operators]\ \\
For $T,S \in \Ls(X,Y)$, we have $T^*,S^* \in \Ls(Y^*, X^*)$, then
\begin{enumerate}[label = (\arabic*)]
    \item $(aT + bS)^* = aT^* + b S^*,\ a,b \in \R \& \C.$
    \item $(a T)^* f(x) = f(a Tx) = af(Tx) = aT^* f(x)$.
    \item $(ST)^* = T^* S^*$.
\end{enumerate}
\end{proposition}
\begin{remark}\ \\
From (2), if $T \in \Ls(X,X)$ is invertible, then $T^* \in \Ls(X^*,X^*)$ is invertible and $(T^*)^{-1} = (T^{-1})^*$.
\end{remark}

\begin{remark}\ \\
Just as Hilbert space, we have the notion of orthogonality.
\end{remark}

\begin{definition}\ \\
Let $A \subset X$ be a Banach space, the orthogonal complement $A^\perp$ of $A$ is a subset of $X^*$ s.t.
\begin{equation*}
    A^\perp = \{f \in X^*:\ f(x) = 0,\ \forall\ x \in A\}.
\end{equation*}
$A^\perp$ is a closed linear subspace of $X^*$.
\end{definition}



\begin{proposition}\ \\
Let $T \in \Ls(X,Y)$, $T^* \in \Ls(Y^*, X^*)$, where $X,Y$ are Banach spaces, then
\begin{center}
    $(\Im{T})^\perp \subset Y^*$ and $\ker{T^*} \subset Y^*$
\end{center}
satisfy $(\Im{T})^* = \ker{T^*}$.
\end{proposition}
\begin{proof}\ \\
$f \in \Im{T}^\perp \Leftrightarrow f(Tx) = 0,\ \forall\ x \in X$, which implies $T^* f = 0$ and thus $f \in \ker{T^*}$.
\end{proof}

\np \textbf{Specialize to Hilbert space}:\\
Suppose $\Hs = \text{ Hilbert space }$, then by \hyperref[RRT]{Riesz Representation Theorem}
\begin{equation*}
    \Hs^* = \Hs,
\end{equation*}
i.e. $f \in \Hs^* \Leftrightarrow\ \exists\ y \in \Hs,\ f(x) = \inprod{x}{y},\ x \in \Hs$.\\
Let $T \in \Ls(\Hs,\Hs),\ T^* \in \Ls(\Hs^*,\Hs^*)$, we have
\begin{equation*}
    T^*f(x) = f(Tx) = \inprod{Tx}{y},\ x \in \Hs,\ f \in \Hs^*.
\end{equation*}
Write $T^* f(x) = \inprod{x}{T^*y}$, this defines $T^*y$ and $T^*: \Hs \to \Hs$.\\
Hence $\inprod{Tx}{y} = \inprod{x}{T^* y},\ x,y \in \Hs$. Clearly, $R^*$ is a bounded linear operator on $\Hs$, i.e, $T^* \in \Ls(\Hs, \Hs)$, and we have $\norm{T^*} = \norm{T}$, since
\begin{equation*}
    \norm{T^*} = \sup_{\norm{y} = 1} \norm{T^* y} = \sup_{\norm{y} = 1,\ \norm{x} = 1} \inprod{x}{T^* y} = \sup_{\norm{y} = 1,\ \norm{x} = 1} \inprod{Tx}{y} = \norm{T}.
\end{equation*}
Therefore, $T^* \in \Ls(\Hs^*,\Hs^*) \Rightarrow T^* \in \Ls(\Hs, \Hs)$ via Riesz.\\
Note that if $T^* \in \Ls(\Hs,\Hs)$, we have
\begin{equation*}
    (aT)^* = \bar{a} T^*,\ a \in \C,\ \bar{a} = \text{ complex conjugate }.
\end{equation*}

\np Note that $T \in \Ls(X,Y) \Rightarrow T^* \in \Ls(Y^*, X^*)$, and $\Im{T}^\perp = \ker{T^*}$. For Hilbert space this implies
\begin{equation*}
    \bar{\Im{T}} \bigoplus \ker(T^*) = \Hs,
\end{equation*}
i.e. orthogonal decomposition. Since $\ker{T^*} = (\Im{T})^\perp$, we use fact that if $E \in \Hs$, $E$ be a subspace,
\begin{equation*}
    (E^\perp)^\perp = \bar{E} = \text{ closure of } E.
\end{equation*}
Apply this to prove the \underline{Von-Newton weak ergodic theorem}.


\np \textbf{Setup:} probability space $(\Omega,\ \Fs,\ \mathbb{P})$, $\mathbb{P} = 1$.\\
Let $T: \Omega \to \Omega$ be a measurable map, we say $T$ is \textit{measure preserving} if
\begin{equation*}
    \P(T^{-1} A) = \P(A),\ A \in \Fs.
\end{equation*}
where $T^{-1}A = \{\omega \in \Omega:\ T \omega \in A\}$. 

\begin{example}\ \\
$\Omega = [0,1]$ Lebesgue measure $\P$, $\Fs = \text{ Borel } \sigma-\text{algebra}$. For $\lambda \in \R$, let
\begin{equation*}
    T \omega = \omega + (\lambda \text{ mod } 1).
\end{equation*}
This is equivalent to rotation on the unit sphere through an angle $2\pi \lambda$. $T$ is measure preserving and one-to-one correspondence. $T^{-1}$ exists.
\end{example}

\begin{example}\ \\
$\Omega = [0,1]$, $\P = \text{ Lebesgue measure },\ \Fs = \text{ Borel sets }$. Let
\begin{equation*}
    T \omega = 2 \omega \text{ mod } 1.
\end{equation*}
$T$ is the shift operator on the binary representation, i.e. $\omega = \summ{j=1}\infty \frac{a_j}{2^j},\ a_j = 0 \text{ or } 1$. Therefore,
\begin{equation*}
    T \omega = \summ{j = 1}\infty \frac{a_{j + 1}}{2^j}.
\end{equation*}
Let
\begin{equation*}
    I_{n,k} = [\frac{k - 1}{2^n}, \frac{k}{2^n}],\ 1 \leq k < 2^n,
\end{equation*}
which is the Dyadic interval. We have
\begin{equation*}
    T^{-1} I_{n,k} = I_{n+1,k} \cup I_{n+1, k+ 2^n}.
\end{equation*}
We have $\P(T^{-1} I_{n,k}) = \P(I_{n,k})$ for all Dyadic interval, which implies
\begin{equation*}
    \P(T^{-1} 0) = P(0),\ \forall\ 0 \in \Fs.
\end{equation*}
Hence $T$ is measure preserving but $T$ is \underline{not} one-to-one correspondence. In fact, $T$ is a $2 \to 1$ mapping. Action of $T$ is
\begin{align*}
    [0,\frac{1}{2}] &\overset{T}{\to} [0,1],\\
    [\frac{1}{2},1] &\overset{T}{\to} [0,1].
\end{align*}
$T$ doubles the length of a dyadic interval, and it is measure preserving since it is $2 \to 1$. $T$ is an expanding map, such mappings are called \textit{hyperbolic}.
\end{example}

\vspace{12pt}
\subsubsection{Von Newmann Ergodic Theorem}

\np Suppose $T: \Omega \to \Omega$ is measure-preserving, we can associate operator $U$ on $\Ls^2(\Omega)$ by $U f(\omega) = f(T \omega),\ \omega \in \Omega$. Then we have
\begin{equation*}
    \int_\Omega f(T \omega) \dr \mu(\omega) = \int_\Omega f(\omega) \dr \mu(\omega),\ \forall\ f \in \Ls^1(\Omega)
\end{equation*}
It is true if $f = \1_A$ and then extend to $\Ls^1(\Omega)$.\\
Suppose $\vphi \in \Ls^2(\Omega)$, we have $U \vphi(\omega) = \vphi (T \omega)$. Also,
\begin{align*}
    \inprod{U \vphi}{U \psi} &= \int_\Omega \vphi(T \omega) \psi(T \omega) \dr \mu(\omega)\\
    &= \int_\Omega \vphi(\omega) \psi(\omega) \dr \mu(\omega)\\
    &= \inprod{\vphi}{\psi}.
\end{align*}
Conclude that
\begin{equation*}
    \inprod{U \vphi}{U \psi} = \inprod{\vphi}{\psi},\ \vphi,\psi \in \Ls^2(\Omega).
\end{equation*}
$U$ is a bounded linear operator on $\Hs = \Ls^2(\Omega)$ with $\norm{U} = 1,\ \norm{U\vphi} = \norm{\vphi},\ \vphi \in \Hs$.\\
In addition,
\begin{equation*}
    \inprod{U \vphi}{U \psi} = \inprod{\vphi}{\psi} \Rightarrow \inprod{U^* U \vphi}{\psi} = \inprod{\vphi}{\psi},\ \vphi,\psi \in \Hs,
\end{equation*}
which implies $U^*U = I$. Therefore, $U$ is one-to-one but not necessarily onto.

\np If $U$ is onto then $UU^* = U^* U = I$ $\Rightarrow$ $U$ is a unitary operator on $\Hs$, hence $U$ is invertible. $U$ is invertible if and only if $T$ is one-to-one. Then $U^* \vphi(\omega) = \vphi(T^{-1} \omega),\ \omega \in \Omega$.
\begin{remark}\ \\
$T$ (1-1) implies $T$ almost onto. To see this let $A$ be a set s.t. $T(\Omega) \subset A$, and hence $T^{-1} A = \Omega,\ \P(T^{-1}A) = \P(\Omega) = \Delta$ $\Rightarrow$ $\P(A) = 1$ $\Rightarrow$ $\P(\Omega \setminus A) = 0$. 
\end{remark}

\np In the case $T$ is not invertible (e.g. 2-1 mapping), might expect similar formula for $U^*$. In the shift operator case, $T_1:[0,\frac{1}{2}] \to [0,1],\ T_2: [\frac{1}{2},1] \to [0,1]$, and $T_1, T_2$ are invertible, we have
\begin{equation*}
    U^* \vphi(\omega) = \frac{1}{2}[\vphi(T_1^{-1} \omega) + \vphi(T_2^{-1} \omega)].
\end{equation*}

\vspace{3pt}
\begin{definition}\ \\
A measure preserving mapping $T: \Omega \to \Omega$ is \textit{ergodic} if and only if the only eigen function $\vphi \in \Ls^2(\Omega)$ of the corresponding operator $U$ is the constant function, i.e.
\begin{equation*}
    U \vphi = \vphi \Rightarrow \vphi \equiv \text{constant}.
\end{equation*}
\end{definition}

\vspace{3pt}
\begin{lemma}\ \\
A measure preserving mapping $T: \Omega \to \Omega$ is \textit{ergodic} if and only if invariant sets of $T$ have probability $0$ or $1$, i.e. if $A \in \Fs$ satisfies 
\begin{equation*}
    \P[(A - T^{-1} A) \cup (T^{-1} A - A)] = 0, 
\end{equation*}
then $\P(A) = 0$ or $\P(A) = 1$.
\end{lemma}
\begin{proof}\ \\
Assume $T$ is not ergodic, then $\exists\ \vphi \in \Ls^2(\Omega),\ U \vphi = \vphi$. Hence can find $a,b \in \R,\ a < b$ s.t. 
\begin{equation*}
    A = \{\omega \in \Omega:\ a < \vphi(\omega) < b\}
\end{equation*}
has $0 < \P(A) < 1$. However,
\begin{align*}
    T^{-1}A &= \{\omega:\ T \omega \in A\}\\
    &= \{\omega:\ a < \vphi(T \omega) < b\}\\
    &= \{\omega:\ a < \vphi(\omega) < b\} = A,
\end{align*}
and thus $A$ is invariant.

\np Conversely suppose $A \in \Fs$, we have $A = T^{-1}A$ up to measure-zero sets and $0 < \P(A) < 1$, then $\vphi = \1_A$ satisfies $U \vphi = \vphi,\ \vphi \in \Ls^2(\Omega),\ \vphi \not\equiv \text{constant}$.
\end{proof}


\vspace{6pt}
\begin{theorem}[Von Newmann Ergodic Theorem, v.1] \label{VN ergodic}\ \\
Suppose $T: \Omega \to \Omega$ is measure preserving, then for any $\vphi \in \Ls^2(\Omega)$, one has 
\begin{equation*}
    \lim_{N \to \infty} \frac{1}{N} \summ{n=0}{N-1} \vphi(T^n \cdot) = \int_\Omega \vphi(\omega) \dr \P(\omega).
\end{equation*}
\end{theorem}
\begin{remark}\ \\
Convergence is in the $\Ls^2(\Omega)$ sense, i.e. mean square.
\end{remark}

\begin{proposition}[Von Newmann, v.2]\ \\
$T: \Omega \to \Omega$ measure preserving, $\vphi \in \Ls^2(\Omega),\ \E[\vphi] = 0$, then
\begin{equation*}
    \frac{1}{N} \summ{n = 0}{N - 1} \vphi(T^n \cdot) \to 0
\end{equation*}
in $\Ls^2(\Omega)$ as $N \to \infty$. 
\end{proposition}
\begin{proof}\ \\
Note it suffices to assume $\E[\vphi] = 0$. WTS 
\begin{equation*}
    \lim_{N \to \infty} \frac{1}{N} [I + U + U^2 + \dots + U^{N-1}] \vphi(\cdot) = 0
\end{equation*}
in $\Ls^2(\Omega)$. If $\vphi$ is orthogonal to the constant constant function. Since $\E[\vphi] = 0$, then $\inprod{\vphi}{1} = 0$. Define a "derivative" operator on $\Ls^2(\Omega)$ s.t.
\begin{equation*}
    D \vphi = (U - I) \vphi = \vphi(T \cdot) - \vphi(\cdot).
\end{equation*}
Use Fundamental Theorem of Calculus argument,
\begin{equation*}
    [I + U + U^2 + \dots + U^{N-1}] D\vphi = (U^N - I) \vphi.
\end{equation*}
Hence,
\begin{equation*}
    \norm{\frac{I + U + U^2 + \dots + U^{N-1}}{N} \vphi} \leq \frac{2 \norm{\psi}}{N}
\end{equation*}
if $\vphi = D \psi$. In that case $\lim$ as $N \to \infty$ is zero, i.e. if $\vphi \in \Im{D} \subset \Hs = \Ls^2(\Omega)$, then finished. Note also that
\begin{equation*}
    \norm{\frac{I + U + U^2 + \dots + U^{N-1}}{N}} \leq 1
\end{equation*}
since $\norm{U} = 1$. Hence converge to zero if $\vphi \in \bar{\Im{D}} = \text{ closure of } \Im{D}$. 
\begin{equation*}
    \vphi \in \bar{\Im{D}} \Rightarrow \exists\ \vphi_\ep \in \Im{D},\ \norm{\vphi_\ep - \vphi} < \ep,
\end{equation*}
which implies $\norm{\frac{I + U + \dots + U^{N-1}}{N} (\vphi_\ep - \vphi)} < \ep$.

\np Recall $\bar{\Im{D}} \bigoplus \ker{D^*} = \Hs = \Ls^2(\Omega)$. It suffices to show $\ker{D^*}$ is spanned by constant functions. 

\np Note $T$ ergodic implies $\ker{D}$ is spanned by constants, we have $D \vphi = 0 \Leftrightarrow U\vphi = \vphi$, and
\begin{equation*}
    (D^* \vphi = 0 \Leftrightarrow U^* \vphi = 0) \Rightarrow (\inprod{\vphi, U^* \vphi}{\vphi} = \inprod{\vphi}{\vphi}). 
\end{equation*}
Therefore,
\begin{align*}
    \inprod{U \vphi}{\vphi} &= \inprod{\vphi}{\vphi}\\
    \int \vphi(T \omega) \vphi(\omega) \dr \P(\omega) &=\int \vphi(\omega)^2 \dr \omega\\
    &= \int \vphi(T \omega)^2 \dr \omega,
\end{align*}
which implies
\begin{equation*}
    \frac{1}{2} \int [\vphi(T \omega)^2 + \vphi(\omega)^2] \dr \P(\omega) - \int \vphi(T \omega) \vphi(\omega) \dr \P(\omega) = 0.
\end{equation*}
i.e. $\frac{1}{2} \int[\vphi(T \omega) - \vphi(\omega)]^2 \dr \P(\omega) = 0$, which means
\begin{equation*}
    \vphi(T \omega) = \vphi(\omega),\ \omega \in \Omega.
\end{equation*}
i.e. $\vphi \equiv \text{constant}$ by ergodicity.
\end{proof}

\vspace{12pt}
\subsection{Open Mapping Theorem}

\np Suppose $T: X \to Y$ bounded linear operator on Banach spaces, and $T$ is injective and surjective, i.e. $T^{-1}: Y \to X$ exists. \underline{Open Mapping Theorem} implies $T^{-1}$ is a bounded operator. The main input into argument is \underline{Baire Category Theorem}. 

\begin{definition}\ \\
A set $S$ in a metric space $M$ is \underline{nowhere dense} (e.g. Cantor set) if its closure $\bar{S}$ has empty interior. 
\end{definition}

















\end{document}
